% Defense preparation notes — personal, NOT part of the thesis build.
% Standalone: latexmk -pdf defense-qa.tex   (run from this directory)
\documentclass[11pt,a4paper]{article}

\usepackage[T1]{fontenc}
\usepackage[utf8]{inputenc}
\usepackage[main=english,brazil]{babel}
\usepackage{amsmath}
\usepackage{amssymb}
\usepackage{siunitx}
\usepackage{booktabs}
\usepackage{microtype}
\usepackage{url}
\usepackage[margin=2.5cm]{geometry}
\usepackage[hidelinks]{hyperref}
\pdfstringdefDisableCommands{% math in question titles must not leak into bookmarks
  \def\times{x}\def\rho{rho}\def\approx{~}\def\leq{<=}\def\geq{>=}%
  \def\emph#1{#1}\def\texttt#1{#1}%
}

% --- Q/A macros -------------------------------------------------------------
% A question heading, followed by the answer I intend to give.
\newcommand{\Q}[1]{\subsection*{Q. #1}\addcontentsline{toc}{subsection}{#1}}
\newcommand{\A}{\par\noindent\textbf{A.}\enspace}
% The objection hiding inside the question, or the thing I must not say.
\newcommand{\Trap}[1]{\par\smallskip\noindent\textit{Trap:}\enspace\textit{#1}}
% The artifact, table, or number to point at when answering.
\newcommand{\Point}[1]{\par\smallskip\noindent\textit{Point to:}\enspace #1}
% A concession I have decided in advance to make, quickly and without defending.
\newcommand{\Concede}[1]{\par\smallskip\noindent\textit{Concede:}\enspace #1}

\title{Defense Q\,\&\,A\\[0.5ex]
{\large Preparation notes}}
\author{Giovanni Oliveira}
\date{\today}

\begin{document}
\maketitle

\section*{What this document is}

These are my personal preparation notes for the defense: a running list of the
questions I expect the committee to ask, paired with the answer I actually
intend to give. It is a working document, not a chapter, and it deliberately
sits outside the thesis build (\texttt{tex/tese.tex}). Nothing here is written
for the reader of the dissertation --- it is written for me, the week before.

The purpose is rehearsal, not exposition. Each entry pairs a question with the
shortest honest answer and, where it matters, with the number, figure, or
artifact I should point to when giving it. Where an answer has a weak spot, I
write the weak spot down here rather than discovering it in the room. Where a
question is really an objection wearing a question's clothes, I note what the
objection is and decide \emph{in advance} whether I concede it, scope it, or
push back. Conceding a real limitation quickly is almost always cheaper than
defending it badly.

Much of this list is not guesswork. The manuscript went through ten rounds of
adversarial review, and every correction those rounds forced is a question a
sharp examiner can ask independently: the headline moved from
$0.27$--$0.31\times$ to $0.62\times$ when an artifact and an actuation-timing
bug were found; the encoder/silicon gap was reattributed from a
$1.4$--$1.6\times$ model error to an execution-engine difference; the
language-pool result was narrowed from ``anticipation'' to ``conversion under a
held cursor.'' I would rather be the person who volunteers each of those than
the person caught by them. The audit history is a strength if I present it as
one, and a wound if I let it be excavated.

Three rules I am writing down so that I follow them under pressure:

\begin{itemize}
  \item \textbf{Lead with the claim, then the evidence.} Answer in one sentence
  before elaborating. If the honest answer is ``I don't know'' or ``that is out
  of scope,'' say that first, then explain why the boundary is where it is.

  \item \textbf{Never oversell.} Every quantitative claim has a measurement
  behind it and a condition under which it fails. I should be able to state both
  without being asked. A committee trusts the candidate who volunteers the
  failure mode.

  \item \textbf{Separate what I measured from what I modeled.} Most of the hard
  questions land exactly on that seam. The answer usually begins by saying which
  side of it a given number came from --- simulated (the week economics),
  measured on silicon (the encoder cost model), or replayed through a real
  control plane with mocked engines (the live validation).
\end{itemize}

% ============================================================================
\section{The story spine: five beats and their numbers}

Everything below is filed under one of five beats, in the order the argument
is built and in the order a defense actually unfolds. If I can recall the
five beats and one number each, I can rebuild the whole thesis in the room
from memory --- and I can place any question the committee asks, because
every question is an attack on exactly one beat.

\begin{enumerate}
  \item \textbf{PREMISE --- the signal is decoupled.} The demand that drives
  the vision pool is not carried by the request stream every autoscaler
  watches. \emph{Killer number:} $R^2 = \num{0.123}$ (images versus requests;
  $0.654$ for decode offered-work under the identical construction).

  \item \textbf{FORECAST --- and it is operationally forecastable.} The same
  signal is the most persistent of the four, so pretrained models read it
  zero-shot at provisioning horizons. \emph{Killer number:} autocorrelation
  half-life \SI{69}{\minute} against \SI{9}{\minute} for requests.

  \item \textbf{CONTROL --- watching it is not enough; placement and
  conversion are.} Better information through the wrong conversion buys
  nothing, and trailing placement loses the fast onset.
  \emph{Killer number:} the burst1@L60 triad,
  1{,}616 / 1{,}717 / \SI{736}{\milli\second} (reactive / forecast through
  mean conversion / forecast through tail-aware conversion).

  \item \textbf{ECONOMICS --- the same compliance costs less.}
  Equal declared compliance at less standing capacity, never a latency win.
  \emph{Killer number:} $0.623\times$ --- 669.7 against 1{,}074.8
  worker-hours over the 5.60-day, 749{,}601-request evaluation span.

  \item \textbf{VALIDITY --- each layer is anchored where it is cheapest.}
  Silicon for the encoder cost, a live control plane for the planner, a
  sensitivity bracket for the H100 timing; no single all-real experiment.
  \emph{Killer number:} $12.8\%$ --- the composed cost model against the
  serving engine's own vision tower over 69 batch sizes, $N \leq 64$.
\end{enumerate}

Then three closing layers, which are not beats but preparation:
\textbf{theory and definitions} (Section~\ref{sec:theory}), so no term is
fumbled; \textbf{scope and concessions} (Section~\ref{sec:scope}), the
limitations I have decided in advance to volunteer; and the \textbf{cram
layer} (Section~\ref{sec:cram}), the numbers that must come out in under
five seconds.

\paragraph{On the hardest questioner.} The statistics and queueing questions
--- filed here into beats~1 and~2 --- come from a committee member with
queueing-theory or time-series background. These reward precision and punish
hand-waving more than any others.

\paragraph{On the kill shots.} Several questions in this document are the
hardest attacks a hostile expert can build, chaining limitations I concede
individually into a headline-level attack. For each I write the wound plainly
and the move I have decided on. Where the honest move is to concede and
re-scope, I do; a defended overclaim there loses the room. Three of them (the
actuation-timing wording, the $4\%$-noise sentence, the two per-stage bounds)
are also manuscript fixes I should make before the defense, not only answers.
\tableofcontents

% ==========================================================================
\section{Opening: the thesis in one breath}

Rehearse these until they are automatic. The opening minutes set whether the
committee reads the rest as a defense or an interrogation.

\Q{Summarize the thesis in two minutes.}
\A A modality-disaggregated multimodal serving system runs its vision encoder
in a pool of its own, separate from the language model's. The question is how to
size each pool. The finding is that no single load signal can do it: in a
production trace, the image demand that drives the encoder is nearly decoupled
from the request rate a conventional autoscaler watches
($R^2 = \num{0.123}$), an order of magnitude burstier, and its extremes do not
even co-occur with the request signal's. So the encoder pool must be
provisioned on its own signal. But a floorless, matched-rating trailing
controller watching that signal still fails, because bursts ramp to peak in
a median of six minutes against replica cold starts of tens of seconds to
minutes (standing reactive floors do pass, at higher cost). The escape is that
the signal most dangerous to reactive scaling is also the most \emph{persistent}
(autocorrelation half-life 69 minutes, Hurst ${\approx}0.93$) --- and therefore
the most forecastable. Pretrained time-series foundation models forecast it
zero-shot, beating every classical baseline in all twenty (signal, resolution)
cells. Driving the encoder schedule from that forecast holds the tail SLA
on all five declared evaluation windows at $0.62\times$ the provisioned
worker-time of the deployable static fleet, where the identically informed
floorless trailing controller loses the fastest onset at every lead time. The finer
daily-fold audit surfaces one further \emph{onset} the pure forecast
misses (fold F6, p99 \SI{1264}{\milli\second}), recovered post hoc by the
hybrid combiner at $0.76\times$; a second fold, F2, fails through
incomplete fixed-context archive coverage, not a missed onset. So the
honest headline is ``all five declared windows at $0.62\times$, one
additional daily-fold onset miss recovered post hoc at $0.76\times$ (plus
one coverage-failure fold),'' not ``every onset.''

\Q{What is the single most important number in the thesis?}
\A $R^2 = \num{0.123}$. Everything else follows from it. If image demand were
predictable from request rate, there would be no provisioning gap, no need for a
per-pool signal, and no thesis --- you would just scale the encoder pool on the
request rate like everything else. The $0.62\times$ is the payoff, but the
$0.123$ is the reason the work exists.
\Trap{Do not answer ``$0.62\times$.'' That invites the whole defense to be
litigated on the economics, which is the part with the most modeling assumptions
and the most audit history. The descriptive core is where I am strongest.}

\Q{What would a practitioner take away tomorrow?}
\A Aggregate the image count your request stream already declares at admission
--- it needs no new instrumentation --- and scale the encoder pool on a forecast
of \emph{that}, with a p90 band for headroom, not on request rate and not
reactively. And convert demand to replicas against a measured tail knee, not a
mean.

% ==========================================================================
\section{Beat 1 --- PREMISE: the signal is decoupled}
\label{sec:beat1}

The premise beat. Every question here attacks the same claim: that the vision pool's demand is not carried by the request stream. It is the beat I am strongest on, because it is measured rather than modeled, and the beat I should steer toward when the economics get litigated. The statistics questions belong here too --- burstiness, persistence, and the estimators behind them are all premise evidence --- as do the prior systems that observed per-modality independence without quantifying it.

\Q{Why does this problem matter --- is it not just ``autoscale on the right metric''?}
\A ``The right metric'' is my \emph{conclusion}, not my premise --- and each of
its three clauses had to be established against a live alternative. First, that
the right metric exists and is not the one everyone watches: image demand
against request rate is $R^2 = \num{0.123}$, and a monitor watching request
rate, mean utilization, or any minute-averaged series sees a \emph{healthy}
pool during an SLA violation --- p99 encoder delay hits \SI{5.0}{\second} at
29\% mean utilization, at the trace's own recorded rate. Second, that watching
the right metric is not enough: a reactive controller reading the true image
signal still loses the fastest-onset burst at every realizable issued
cold-start lead from 15 to \SI{165}{\second} (the old \SI{300}{\second}
cell was target-indexed, an evaluation bound rather than an issued lead;
p99 encoder delay 9.5--\SI{54}{\second},
degrading toward the \SI{54.2}{\second} never-scaled limit). Third, that the
escape --- forecasting --- is
actually open, which nothing guarantees a priori: burstiness and
forecastability could have anticorrelated. They do not; the burstiest signal
carries the longest memory. Each clause is a measured result, and reactive
autoscaling on aggregate signals is not a strawman --- it is the shipped
default in the platform I build on and the inherited practice from stateless
web serving.

\Q{Is this not obvious? Of course an image pool scales with images.}
\A The direction is obvious; the three quantities that make it a provisioning
problem are not, and prior work got none of them. (a)~The \emph{magnitude}:
$R^2 = 0.123$ is not ``weakly coupled,'' it is near-total decoupling ---
the week's top-percentile minutes of the two signals share 5 of ${\sim}101$
members, Jaccard 0.02, because bursts are 93--96\% pure frames-class floods
that the request stream barely registers ($370$\,img/s at under 4.5\,req/s).
(b)~The \emph{insufficiency of the obvious fix}: scale on images reactively
and you still fail --- onsets ramp in a median of six minutes, the fastest in
one, with a median steepest single-minute jump of 208\,img/s, against replica
leads of tens of seconds to minutes. I ran exactly that controller --- the
floorless trailing arm at matched rating --- and it loses
the fast onset at every lead. (A standing reactive \emph{floor} deep enough
does pass; every tested passing floor costs more than the fixed forecast
point on this trace.) (c)~The \emph{escape hatch}: the same signal has
a 69-minute autocorrelation half-life and Hurst ${\approx}0.93$, so the
information a reactive loop lacks is recoverable by a forecaster. ``Obvious''
predicts (a); it does not predict (b) or (c), and (b) is where deployed
systems actually die. ModServe --- the system that introduced this
architecture --- scales its image pool \emph{reactively}; if the full answer
were obvious, the system defining the field would not have shipped the failing
half of it.
\Point{The reactive-on-the-same-signal ablation --- it isolates the
\emph{placement of standing probabilistic headroom}, not signal choice, as
the operative advantage. Say ``placement,'' not ``anticipation'': on burst3
there is no usable anticipation (that is measured on the gradual burst4
onset), and the ablation matches rating and mechanics, not worker-time.
Without it the thesis is a metric-selection paper; with it, the forward band
is load-bearing.}

\Q{Is $R^2{=}0.123$ not trivial --- of course request rate is not the same thing as image content?}
\A The triviality is in a claim I do not make. RQ1 asks whether the
\emph{rate} signal --- the input conventional autoscalers, and every
disaggregated-serving autoscaler in the literature, actually watch ---
carries the vision pool's demand. It does not ($R^2{=}0.123$), while it
\emph{does} carry decode ($R^2{=}0.654$) under the identical
arrival-stamping construction, and the two signals' extreme minutes are
near-disjoint (Jaccard 0.02). That contrast is the non-trivial part: the
same conventional trigger works for one pool and fails silently for the
other (p99 \SI{5.0}{\second} at 29\% mean utilization). And the image
count \emph{is} knowable --- every request declares it at admission --- so
this is a correctable choice of \emph{which} signal to watch, not an
information-theoretic impossibility.
\Trap{Do not drift into ``image demand is unknowable.'' That contradicts my
own text (the signal is declared at admission and is the \emph{most}
forecastable). The claim is: knowable, but not from the signal deployed
autoscalers watch.}

\Q{$R^2{=}0.123$ is a linear, through-origin fit. Where is the confidence interval, and did you test nonlinear or rank coupling --- Spearman, mutual information?}
\A The decoupling was tested under three relationships, not one:
constant-fraction scaling ($\alpha{=}15.1$ img/req, forced through origin),
rank correlation, and lagged linear cross-correlation. Supporting figures:
Pearson $r{=}0.36$ images-vs-rate against 0.81 for decode; at request level
image-count-vs-generation Spearman $-0.06$, with a non-monotone per-bucket
median; window-level cross-correlation peaking at lag zero with no lead. I
do not overclaim: nonlinear or multivariate coupling is ``not ruled out by
these tests,'' and I defer it to the covariate experiment --- a stronger,
model-based probe than mutual information --- which comes up empty (the
request covariate explains only $R^2{=}0.12$ of images; covariate effect
within $\pm0.4\%$ at provisioning cadence).
\Concede{No confidence interval on 0.123, and no mutual-information test.
Concede both; argue the zero-shot covariate experiment already tests
exploitable nonlinear coupling and finds none.}

\Q{Why does averaging not rescue the reactive loop? Utilization was only 29\%.}
\A Because the arrival process is overdispersed at \emph{every} operational
scale: the Fano factor of the counting process grows from 1.8 at one second
to 938 at one hour, instead of flattening as it would for anything
Poisson-like ($\text{Fano}=1$ at all scales). There is no averaging window at
which the signal becomes safe for a threshold rule --- widen the window and
the variance grows faster than the mean. That is the operational meaning of
long-range dependence, and it is the same pathology Leland et al.\ proved for
Ethernet in 1994: burstiness that does not smooth under aggregation. The 29\%
figure is the sharpest form of the point --- the pool violates its p99 bound
at 29\% \emph{mean} utilization because episodic frames-class floods induce
head-of-line blocking behind the FIFO while every minute-averaged dashboard
series stays green. Mean utilization is the wrong sufficient statistic for a
tail SLA under an LRD arrival process; that is a structural claim, not a
tuning complaint.

\Q{``Episodic head-of-line blocking at 29\% mean utilization'' --- is that a queueing result, or a deterministic replay of one known spike in queueing vocabulary? There is no arrival model, no service-time distribution, no $\rho$, no M/G/1.}
\A It is a deterministic event-driven FIFO replay, and I will not dignify it
as an analytic queueing model --- no ``G,'' no Pollaczek--Khinchine term, no
stability region, and I claim none. The result is a systems observation: a
FIFO of $C$ finite workers, fed a second that lands 2{,}726 images, backs up
past 1{,}000 queued images for 805 seconds and drains between excursions, so
the \emph{mean} utilization (29\%) is a near-useless summary of a bimodal
occupancy while p99 hits \SI{5.0}{\second}. ``Head-of-line blocking'' is
used in its correct systems sense, and I show it is not a dispatch-cap
artifact (p99 5.2/5.0/\SI{6.7}{\second} at caps 32/64/128).
\Trap{Never invoke M/G/1, PASTA, or Little to dignify it --- that invites a
demand for the service-time distribution and stability analysis I do not
have. Call it a deterministic replay exhibiting tail blocking a mean monitor
misses.}

\Q{How did you estimate Hurst ${\approx}0.93$? Did R/S, DFA, and a spectral estimator agree, or is this one number dressed up as evidence of long-range dependence?}
\A One estimator, and the thesis prints that caveat verbatim: $H$ is derived
algebraically from the spectral slope, $H=(1+\beta)/2$ --- ``one estimator
viewed two ways, not two corroborating measurements.'' No R/S, DFA, or
wavelet estimator was run. Two honesties follow. The slope is fit only over
the sub-diurnal band (10 minutes to 12 hours), clear of the diurnal
fundamental, so the LRD claim is deseasonalized by band selection; the
raw-series ACF half-life (69 minutes) is \emph{not} deseasonalized and I use
it only comparatively. And I decline any joint claim where the estimators
disagree --- the decode series pairs the shortest half-life with the
second-highest $H$, and I say so. Everything load-bearing rests on the
images-versus-requests \emph{contrast} ($69$ vs \SI{9}{\minute}; $H\,0.93$
vs $0.83$), which both signals' shared diurnal component cancels out of.
\Trap{Do not defend $0.93$ as a precise or corroborated LRD constant.
Concede it is one spectral estimator; pivot to the ratio, and to the
model-free ACF half-life and Fano growth, which need no clean self-similar
model.}

\Q{Where are your stationarity tests? A Hurst above $0.5$ and a Fano factor climbing to 938 are statistically confusable with simple non-stationarity on a seven-day record.}
\A No formal ADF/KPSS test is reported, and I concede that LRD and
non-stationarity are not cleanly separable on seven days --- I say in print
that periods of 84 hours and up are record-length artifacts. The
stationarity language is descriptive, and the load-bearing claim is
deliberately weaker than a self-similarity model: the counting process is
overdispersed at \emph{every} aggregation scale an autoscaler could use
(Fano 1.8 at \SI{1}{\second} to 938 at one hour), which is directly measured
within the window and is all a threshold rule needs to be unsafe. The
band-limited spectral fit sidesteps the lowest frequencies where a trend
would dominate.
\Trap{Concede the seven-day identifiability limit rather than defend
asymptotic LRD. The operational claim needs only \emph{within-window}
overdispersion, which is measured, not modeled.}

\Q{Every burstiness number --- Hurst, Fano, the 69-minute half-life, the 18 onsets, the burst p99 --- is a point estimate from one 7-day path with no confidence interval. And your own $H{\approx}0.93$ implies a sample mean's variance falls as $n^{2H-2}$, so 10{,}080 one-minute bins carry an effective sample size of about \emph{four}. On what basis is any of this distributional?}
\A The honest answer is that these are within-path \emph{descriptions}, and I
pre-hedge each: Hurst is ``one estimator viewed two ways, not two
corroborating measurements''; a seven-day record ``cannot establish
asymptotically unbounded'' Fano growth; periods above 84 hours are
record-length artifacts. The claims consumed downstream are \emph{contrasts
within the same realization} --- images 69\,min vs requests 9\,min, CV$^2$
3.83 vs 0.32 --- robust to the shared diurnal component precisely because
both signals ride it, and the generality argument runs through a stated
\emph{mechanism} (batch-video tenants sharing a chat cluster) with an
explicit falsification condition, not through sampling.
\Concede{There is no CI, bootstrap, or ensemble on any burstiness metric.
Concede it and name a block-bootstrap of CV$^2$/half-life as the right
addition. ``Mechanism, not sampling'' rescues the \emph{direction} of a
contrast, never a point estimate's precision --- the effective-$n$ point is
fair.}

\Q{Your ``burstiest signal is the most forecastable'' is sold as a discovered coupling. Is it not just the heavy-tail-to-long-range-dependence identity of one batch-arrival process?}
\A Partly, and I will name the mechanism rather than defend a coincidence.
The frames class (2.7\% of requests, 81.2\% of image volume) drives
\emph{both} the high CV$^2$ (3.83) and the low-frequency spectral power
(Hurst ${\approx}0.93$) --- they are the second moment and the spectral
slope of one heavy-tailed component, and a heavy-tailed batch/ON--OFF
superposition mechanically produces both (the Taqqu--Willinger duality). So
this is one identifiable mechanism, not two properties that surprisingly
coincide. What the operational claim needs is not that they are independent
--- it is that the envelope is \emph{forecastable} at provisioning leads,
and that is measured directly (one-minute-bin WAPE flat within
${\approx}2$ points over twelve steps, p90 coverage 0.88--0.90), not
inferred from the coupling.
\Concede{Downgrade ``surprising coupling'' to ``the operational readout of
one identifiable mechanism.'' Forecastability is validated independently;
do not defend it as a general law of bursty signals.}

\Q{The 18 episodes, the six evaluation windows, and ``which one is burst3'' all come from a single unswept recipe --- image rate above its p99, contiguous runs merged to median crossings. At p98 or p99.5, do you still get 18 episodes, and is burst3 still the worst?}
\A The threshold is principled --- it is the SLA percentile the pools are
graded on --- and the episode \emph{set} is not knife-edge: the frames class
supplies 93--96\% of every top episode's volume and clustering isolates a
single mixed outlier, so the trace contains ``one burst type'' and moving
the cut rescales magnitude and duration, not kind. The load-bearing windows
are the largest episodes, far above any p98/p99 boundary, and the burstiness
characterization (Fano, CV$^2$, ACF) does not use the segmentation at all.
\Concede{I did not sweep p98/p99.5 or the merge gap, and the fastest-onset
statistics --- which power the onset-vs-cold-start argument --- are exactly
the ones most sensitive to where an episode starts. A threshold-sensitivity
table is the right addition; concede it rather than claim untested
robustness.}

\Q{Your entire vision-pool result rides the frames class --- video-like grids, 2.7\% of requests. Ordinary image-chat does not burst or decouple. Strip frames and your decoupling, inversion, and 0.62$\times$ evaporate. Is this a ``multimodal serving'' thesis or a ``video-frame-batch provisioning'' note?}
\A I accept the regime framing, and it is on the record: my own trace chapter
states that a workload where image volume tracks request rate ``would
collapse the decoupling, and with it the case for a separate vision-pool
signal.'' The frames class is not an artifact to strip --- it is \emph{why
modality disaggregation exists}: a single $448^2$ image is cheap, and the
case for a separately scaled encoder pool is created precisely by heavy,
bursty visual work. Removing frames removes the motivation for the
architecture the whole field (ModServe, EPD, ElasticMM, vLLM-Omni) already
ships; it does not refute the thesis, it deletes its regime.
\Concede{The magnitudes ($R^2$, $H$, 0.62$\times$) are this frames-heavy
trace's, and I say so; a frames-light workload weakens the case toward the
language pool's break-even. This is the narrowest empirical flank, and a
second trace is the honest test, not a wider title.}

\Q{Why one trace? Does this generalize beyond Azure's workload?}
\A One trace --- I will say that before anyone asks, and the scope section
says it in print. Three things bound the damage. First, there is no
alternative: AzureLLMInference 2025 is, to my knowledge, the only public
production multimodal trace carrying per-request image counts; the
characterization became possible the day it was released. Second, the
phenomena are not exotic to this trace: heavy-tailed, long-range-dependent
arrivals with a dominant burst class are exactly what Leland et
al.~(1994) measured in Ethernet traffic and Crovella \& Bestavros (1997) in
the Web --- self-similarity is the \emph{normal} finding for aggregated human
demand, and a serving trace showing Hurst $0.93$ is continuous with thirty
years of teletraffic measurement, not an outlier needing luck to reproduce.
Third, the deliverable that generalizes is the \emph{method}: characterize
per-pool signals, test forecastability at provisioning horizons, size against
a measured tail knee. The constants ($R^2$, the knees, $0.62\times$) are this
trace's; the procedure is not, and every step of it is scripted and
re-runnable on the next trace.
\Concede{One trace, a by-construction 50/50 modality split, image counts but
not sizes. All three, immediately, unprompted.}

\Q{Your split is a by-construction 50/50 text/image mix. How much of the result is an artifact of that ratio?}
\A The exact numbers are ratio-sensitive; the mechanism is not, and I keep
them apart. $R^2{=}0.123$ and the $4\times$/$2\times$ static missizing are
this trace's constants --- a heavier image fraction would inflate vision
load and \emph{sharpen} the decoupling, a lighter one shrink the pool. What
does not move is the driver: each burst is a 93--96\% frames-class flood,
and frames are 2.7\% of requests but 81.2\% of image volume, so the
encoder's episodic head-of-line blocking is a within-image-traffic
phenomenon a different text fraction rescales but does not remove.
\Concede{The 50/50 is a dataset property, not a measured field ratio, and
$R^2$ is sensitive to it. I claim the sign and mechanism are
split-invariant, not the constant --- and I say the split is engineered, in
scope.}

\Q{Who is the audience --- systems or ML?}
\A Systems, using ML as an instrument. The research questions are capacity
questions --- what signal sizes a pool, when must capacity exist, what does it
cost --- and the evaluation currency is systems currency: tail SLA compliance
$\times$ provisioned worker-time. The forecasters are consumed as
off-the-shelf, zero-shot components; I benchmark them the way a systems
person benchmarks a disk, and one of the findings is precisely that the ML
choice is second-order (all three foundation models pass the week at the
identical \SI{651.1}{\milli\second} tail; only cost shifts). The thesis makes
no contribution to forecasting models and claims none; it contributes the
problem framing, the measured system, and the control loop.

\Q{ModServe already noted per-modality independence. What is left for you?}
\A ModServe observes it qualitatively and scales the image pool \emph{reactively}
on observed token load; it also reports a 79\% prediction error and a
25--41.3\% serving-cost saving. My claim is narrower and, I think, still real: I
\emph{quantify} the decoupling, frame the per-pool signals as \emph{forecasting
targets} --- benchmarked zero-shot at provisioning horizons --- and consume them
in a closed-loop per-pool scaler. The distinction that survives scrutiny is that
no prior system drives its autoscaler with a \emph{future-horizon forecast} of a
per-pool signal. I narrowed this claim twice under review and it is now stated as
carefully as I know how.
\Trap{An early draft said ``first forecast-framed characterization'' and ``none
asks.'' Both were too strong --- ElasticMM has a proactive step, and
TokenScale/HeteroScale use predictive metrics. Concede those immediately and hold
only the narrow line.}

\Q{ElasticMM is both modality-disaggregated and proactive. What is left for you?}
\A ElasticMM's proactive step is not a demand forecast --- it is a static
allocation of idle instances by a long-run peak-to-average burst-tolerance
heuristic over \emph{aggregate} load, and its own paper calls short-term
per-stream arrivals ``difficult to predict.'' That is exactly the claim I
overturn: the image signal is the \emph{most} forecastable of the four
(Hurst ${\approx}0.93$, 69-minute half-life), zero-shot, at provisioning
horizons. The thesis stands in ElasticMM's declared blind spot --- it
forecasts the per-pool signal ElasticMM treats as unforecastable --- and RQ4
shows anticipatory headroom \emph{timed to the onset} beats a fixed idle
buffer.

\Q{TokenScale and HeteroScale are per-pool, disaggregated autoscalers watching per-stage signals. How is ``per-pool demand signal'' novel against them?}
\A Both are text-only two-stage systems whose ``per-stage'' signals are the
same token stream split by phase. However finely you split prefill from
decode, the vision pool's demand never appears --- it originates in a
distinct request property (image count, $k\times k$ grids) that is
$R^2{=}0.123$-decoupled from the rate. The natural rejoinder, ``add
image-tokens/s as a third velocity,'' is essentially what ModServe does
\emph{reactively}, and the reactive version fails invisibly at the trace's
own rate; the separation is that I \emph{forecast} that signal.

\Q{EPD disaggregation and HydraInfer already split encode, prefill, and decode onto dedicated resources. What is left?}
\A They contribute the placement and scheduling of a three-stage split;
none asks the autoscaling question. EPD and HydraInfer optimize how
encode/prefill/decode share resources and how requests schedule across them,
\emph{within} whatever fleet exists. My contribution is upstream of that:
which signal sizes each pool over time, that the vision pool's signal is
decoupled from request rate and forecastable, and a closed loop that
provisions on it. Their split is my premise; the per-pool predictive
provisioning is the addition.

\Q{RServe overlaps encoding with prefill in a single instance. Your seam serializes them. Does RServe's design not refute your architecture?}
\A RServe occupies the opposite design point deliberately; it is a
different regime, not a refutation. My premise is modality
\emph{disaggregation} --- the encoder as its own independently scaled pool,
the pattern ModServe introduced and ElasticMM and vLLM-Omni pursue --- and
the provisioning question only exists once the pools are separate. Within
that premise my serialized seam (no speculative encoder/prefill overlap) is
\emph{conservative}: overlap could only lower composed TTFT, so it makes my
saturation knees pessimistic, which strengthens an infeasibility claim.
RServe trades the independent scaling I study for latency, and I gesture at
the reconciliation as future work --- the seam is an exact short-horizon
preview of language-pool arrivals, usable for scheduling if not for replica
lead.

% ==========================================================================
\section{Beat 2 --- FORECAST: and it is operationally forecastable}
\label{sec:beat2}

The feasibility beat: nothing guaranteed a priori that the burstiest signal would also be the most forecastable, and if it were not, the escape route closes. Questions here attack the benchmark, the protocol, the metric, and the calibration. The recurring move is to state the unflattering number first, correctly scoped, and then say what the controller actually consumes: not a point forecast, a calibrated upper quantile.

\Q{Are the foundation models actually good, or just better than weak baselines?}
\A Both claims are in the manuscript with their sizes attached, and I will
state the unflattering one first, at the resolution it actually belongs to.
At \emph{one-second} bins the best model's image WAPE is 98.8\% (Toto)
against EWMA's 146.1\%; read against WAPE's no-skill floor (always predicting
zero scores 100\%), that point forecast is barely above no-skill. At the
\emph{one-minute} provisioning-reporting cadence the picture is different:
52.9\% WAPE (Chronos with covariates) against EWMA's 54.2\% --- material
skill against the no-skill floor, on a thin 1.3-point margin over the cheap
baseline. Neither resolution changes what the thesis deploys, because it does
not deploy point forecasts at all: the operating instrument is the
\emph{quantile band}, whose p90 edge is empirically calibrated at provisioning
cadence --- measured coverage 0.88--0.90 against the nominal 0.90, across all
models and signals at one minute. The honest summary: at second resolution
nothing usefully predicts the next bin of a signal whose own CV$^2$ is 3.83;
at minute resolution the point forecast has real but modest skill; and at
either resolution what the models supply calibratedly is how bad it might get
--- which is exactly what provisioning consumes. Wins over the classical
baselines hold in all twenty (signal, resolution) cells regardless, and
degradation across the 12-step horizon at \emph{one-minute} bins is
${\approx}2$ WAPE points, covering the median six-minute onset with margin.
\Trap{Do not quote ``beats every baseline in all twenty cells'' without the
no-skill caveat in the same breath --- and keep that caveat attached to the
\emph{one-second} cell. Calling 98.8\% a one-minute number is the error a
time-series person on the committee will catch; saying it first, correctly
scoped, converts the weakness into credibility.}

\Q{``Beats every baseline in all twenty cells'' has no confidence intervals and no significance test, on 167 burst-skewed one-minute origins, and you concede the three models tie and ARIMA ties Chronos. Is the benchmark not near-degenerate at exactly the cadence you provision on?}
\A ``All twenty cells'' is a \emph{sign} result --- 20/20 same-direction over
independent (signal, resolution) cells --- not a magnitude claim, and I
should have run a Diebold--Mariano across origins rather than bold
sub-0.1-point wins. But the degeneracy is the \emph{premise}, not a bug: the
benchmark's job is to show the shipped predictors do \emph{not} forecast this
signal better (ARIMA 0.668 vs 0.637 at 15\,s; all pass together), so that
when the closed loop separates pass from fail, the separator cannot be
accuracy --- it must be the signal, the quantile, and the conversion. A
near-degenerate accuracy benchmark is what makes ``model choice is
second-order'' a finding.
\Concede{No CIs, no paired test, no multiple-comparison correction, and the
one-minute column is burst-skewed ($n{=}167$). The correct statement at
provisioning cadence is ``not beaten by any baseline, by margins too small
to certify per-cell.''}

\Q{You headline ``foundation models beat every classical baseline in all twenty cells,'' but the classical set is naive, EWMA, and Holt--Winters. Where is ARIMA, the production incumbent?}
\A ARIMA is absent from the accuracy grid and I state that in the
limitations; it is present where it decides an outcome. The reactive
baseline of the results chapters is Dynamo's planner running its ARIMA
predictor \emph{unmodified}, planner-in-the-loop, and the framework-diff
measurement puts Dynamo's shipped predictors on the image signal directly:
ARIMA at \SI{15}{\second} sits at 0.668 WAPE against the foundation models'
0.637 and \emph{ties} Chronos, after which the integer-replica conversion
quantizes the difference away (verdicts unchanged swapping among them). So
the honest scope is: ``beats every classical baseline'' means the three
fixed-parameter families; ARIMA is benchmarked separately and is
second-order on both accuracy and the closed-loop verdict.
\Concede{A full-history, per-series-tuned ARIMA/Holt--Winters under the
identical rolling-origin protocol is the cheapest reviewer-requested
addition, and it would narrow the margins. Offer it; do not pretend it is
done.}

\Q{Why these three foundation models, and not others?}
\A TimesFM, Chronos, and Toto span the design space that existed at
evaluation time --- Google, Amazon, and Datadog lineages; 120M to 2.5B
parameters; different tokenizations --- against naive persistence, EWMA, and
Holt--Winters as the classical frontier, plus Dynamo's four shipped
predictors (ARIMA, Kalman, Prophet, constant) run through the identical
protocol. And the answer to ``why not a \emph{fourth} TSFM'' is that within
the three tested the choice demonstrably does not matter: all three
foundation models pass the week at the \emph{identical}
\SI{651.1}{\milli\second} tail; only provisioning cost shifts (Chronos 518.6
$<$ Toto 625.3 $<$ TimesFM \SI{641.1}{GPU\hour}), because integer-replica
conversion quantizes most of the accuracy difference away. I will not stretch
that into an open-ended ``model X would not matter'' either: the forward
band's \emph{family and shape} are load-bearing --- on burst3 the tested EWMA
calibrations split on outcome (additive and width-scaled fail at p99
$1{,}285$ and $2{,}084$\,ms; multiplicative passes only by carrying
substantially more standing work) --- so a predictor that changes the upper
edge can change the SLA verdict, not just the cost. Within the TSFM class,
identity is second-order after capacity quantization
--- that is a finding, not an evasion --- and the headline deliberately rests
on the \emph{smallest and cheapest} of the three (Chronos, 120M parameters,
versus Toto's 2.5B), so no part of the result depends on forecaster scale.
\Point{Dynamo's own shipped predictors close the loop on this: ARIMA at 15\,s
sits at 0.668 WAPE versus the foundation models' 0.637, Kalman (0.811) trails
even naive persistence (0.764) --- so a forecast-capable planner gains nothing
from predictor choice \emph{within} its shipped set, and the TSFMs are the
cheapest route to the frontier.}

\Q{Your marquee one-minute WAPE numbers are computed on 167 origins covering only the final 20\% of the week --- which happens to contain the largest burst. Is that a representative population?}
\A No, and cross-resolution column comparison is invalid for exactly that
reason --- I flag it in the protocol. The common-span origin threshold is an
index, not wall-clock, so it covers 98.7\% of the week at one-second bins
but only the final \SI{33.4}{\hour} at one-minute bins. The valid
comparisons are model-versus-model \emph{within} a column, which are
unaffected. And the burst over-representation cuts \emph{for} the
provisioning story, not against it: the one-minute models are judged on the
hardest windows, so the reported accuracy is conservative for the use I put
it to.
\Trap{Do not let an examiner bait me into comparing the one-minute and
one-second WAPE columns as if they measure the same population --- they do
not. Same-resolution, model-vs-model only.}

\Q{Chronos-2 and Toto-2 were released after this 2025 trace. How do you know they did not see it in pretraining? Zero-shot collapses if they did.}
\A They could have, and I cannot rule it in or out from the public model
cards --- but the result is protected by a natural control. TimesFM-2.0
\emph{predates} the trace's release entirely and lands within ${\approx}2$
WAPE points of the other two at the one-minute provisioning cadence where
the forecasts are consumed, which bounds any contamination advantage at
about that size. The larger inter-model gaps live in the sub-minute
shot-noise regime, where no model has skill to contaminate.
\Point{TimesFM is effectively a pre-registration control here. Stress that
the headline provisioning cadence is contamination-bounded regardless of
what the two newer models saw.}

\Q{You sell a probabilistic forecast --- nine quantiles --- then grade it with WAPE, a point metric. Why not a proper scoring rule (pinball, CRPS)? And why WAPE over MASE?}
\A Deliberate separation of concerns, and I concede the proper score is not
computed. WAPE is the point-accuracy metric because it is scale-free and
defined on zero-signal bins, where MASE's in-sample-naive denominator is
awkward on a heavily zero-inflated, non-stationary series. Probabilistic
quality is assessed \emph{separately}, by empirical coverage of the specific
edge the scaler consumes: p90 coverage 0.88--0.90 against nominal 0.90 at
one minute, plus a direct exceedance check (11\% of bins pierce the p90 edge
versus a nominal 10\%) on the largest burst. The justification is
fit-for-purpose: the controller reads a single quantile \emph{path}, not the
full distribution, so point accuracy of that path plus calibration of that
one edge is exactly what matters.
\Concede{CRPS/pinball is the textbook choice and is not done. Note too that
the lower band is admittedly broken by zero ties (p10--p90 covers 69\% at
one minute), so a CRPS would be dragged down by the edge I already disown
--- which is why I provision on the upper edge only.}

\Q{Your p90 coverage of 0.88--0.90 is pooled over all twelve horizon steps and over a burst-heavy span. Is the band calibrated at the lead that actually matters?}
\A It is resolved by horizon step, not only pooled. The pooled figure is
0.88--0.90 (Chronos+cov on images 0.888), and beneath it, coverage of the
\emph{upper} edge --- the only one provisioning consumes --- is reported for
every one of the twelve horizon steps on the image signal: $0.876$--$0.903$
at \SI{15}{\second} bins and
$0.857$--$0.911$ at \SI{1}{\minute} bins, with no trend in $h$, so no decay
with lead. Those twelve steps are forecast-evaluation \emph{target horizons}
--- \SIrange{15}{180}{\second} ahead of the forecast origin --- not issued
lead: an $H{=}12$ archive at \SI{15}{\second} bins supplies at most
$(12{-}1)\times15 = \SI{165}{\second}$ of usable issue-time lead, and the
closed loop tests issued leads of 15/30/45/165\,s. Pooled in-burst versus outside-burst coverage is also reported
($0.890$--$0.901$ in burst against $0.890$--$0.896$ outside). The median edge
is 0.43--0.48 and the lower edge is broken by zero-inflation, so I explicitly
do not provision on the full band. The independent cross-check is the
exceedance rate on the largest burst episode itself: 11\% of bins pierce the
p90 against a nominal 10\%.
\Concede{What is \emph{not} measured is onset-stratified, per-episode
calibration --- coverage in the minutes around each of the 18 onsets. Concede
that specific gap; do not concede horizon-step or pooled burst coverage,
which are in the manuscript.}

\Q{Your closed-loop mechanism is p90 headroom standing capacity up \emph{at onsets}. But your calibration --- coverage 0.88--0.90 --- is pooled over the horizon and the span; onset-conditional p90 coverage, the one number the mechanism depends on, is never computed.}
\A Correct about \emph{onsets}, and only about onsets. Horizon-conditional
coverage is measured --- $0.876$--$0.903$ across all twelve steps at
\SI{15}{\second} bins, no trend in $h$ --- and so is burst-conditional
coverage pooled over the models ($0.890$--$0.901$ in burst against
$0.890$--$0.896$ outside), so the band does not decay with lead and does not
collapse inside bursts. Both are upper-edge only; the lower edge is broken by
zero ties, which is why I provision one-sided. Beyond that I have an in-situ
check on the largest burst itself: 11\% of bins pierce the p90 against a
nominal 10\%, and zero calm-day false alarms. That is onset-region evidence,
but a single-episode operational proxy, not a coverage statistic stratified
by onset.
\Concede{The remaining gap is precisely onset-stratified, per-episode
calibration: the mechanism assumes the upper edge stays calibrated
\emph{at} onsets, and I have shown that once, not distributionally. A
per-episode coverage table is the honest addition, and Chronos's occasional
crossed quantiles mean I should say ``a calibrated upper quantile,'' not ``a
calibrated predictive distribution.''}

\Q{Your quantile band ``adds headroom,'' but the point forecast is the median for Chronos and Toto and the \emph{mean} for TimesFM. Is the quantile finding not just an artifact of inconsistent point heads?}
\A It is exactly that, and the thesis diagnoses it rather than being caught
by it. A quantile beats the point in 51 of 160 cells, and the wins
concentrate entirely in models whose point is \emph{not} the median:
TimesFM's mean head is beaten in 18/20 cells (modally by p40), the Gaussian
baselines in 31 more, while for the median-point models quantiles improve
essentially nothing (2 of 80 cells). The mechanism is stated: for a
right-skewed bursty rate the conditional median sits below the mean, so
mean-like forecasts overshoot after spikes and a sub-median quantile
corrects them --- under a common p50 head, TimesFM's lone win reverts to
Toto. The clean operational conclusion: with a median-point foundation model
the point path is already the best single operating line, and the band's
role is headroom, not point correction.
\Point{This is my own finding, not an embarrassment --- lead with the
median-vs-mean-head diagnosis. Disclose too that Chronos-2 emits occasional
locally crossed quantiles, recorded unrepaired apart from a nonnegativity
floor.}

\Q{Why zero-shot forecasting? Would fine-tuning not do better?}
\A Almost certainly, which is exactly why I did not do it. Zero-shot is the
\emph{conservative} reading of the foundation models' capability: it is the lower
bound on what an operator gets on day one, with no trace history and no training
pipeline. If the claim holds zero-shot, fine-tuning can only help. Reporting
fine-tuned numbers would have made the result stronger on paper and weaker as
evidence.

\Q{Why rolling-origin evaluation, and why does the protocol matter?}
\A Because an autoscaler does not consume a forecast the way a benchmark does.
It sits at a point in time, asks for the next $h$ steps, acts, and then the clock
advances. The rolling-origin protocol scores the models exactly that way --- the
forecasts are stitched to tile the timeline as a scaler would consume them. A
static train/test split would have measured something a provisioning system never
experiences.

\Q{The analysis is at one-minute resolution but the controller runs on 15-second bins. Why two cadences?}
\A They serve different jobs. One minute is the provisioning-\emph{reporting}
cadence --- the horizon over which the persistence result and the WAPE
tables are stated, and the scale of burst onsets (median six minutes).
Fifteen seconds is the \emph{actuation} cadence: it matches the granularity
at which a scaler issues decisions and cold-start leads are counted, and the
archive's twelve-step horizon at 15-second bins is what supplies the
up-to-\SI{165}{\second} of issued forecast lead the closed loop consumes.
The forecasters are benchmarked at both (and three more) resolutions, so
nothing is cherry-picked --- the foundation models win all twenty cells
across all five.

\Q{Context tokens are largely image-driven. Did you try combining signals --- forecasting prefill demand as text-borne plus image-driven components?}
\A I did, and it failed for an instructive reason, so I report it as a
closed negative result. Regressing realized context tokens on requests and
images recovers $R^2{=}0.94$ in-sample, the image term carrying 34--36\% of
context mass --- so an additive forecast (text-borne demand plus the scaled
image forecast) looks promising, especially at onsets where prefill regimes
shift. Pooled, it ties or slightly trails the direct context forecast
(within 1.3 WAPE points), and at the onsets there is no robust edge --- the
sharpest detector even flips sign by resolution. The diagnosis is exact:
\emph{at} the onsets the image forecast is the \emph{worst} component
(75--100\% WAPE against 58--72\% for direct context and 33--38\% for
requests), because the image spike is least forecastable precisely where
the structural bet needed it. The direct forecast already absorbs the
image-driven component; structure buys nothing the envelope did not.
\Point{A volunteered negative result: I tested the obvious ``combine the
signals'' idea and it lost to the simpler forecast --- consistent with
covariates being second-order at provisioning cadence.}

\Q{A GPU-resident 120M--2.5B-parameter foundation model versus EWMA, a few lines of arithmetic --- is the accuracy margin worth that operational asymmetry?}
\A At the decision layer the compute is trivial: one 12-step forecast costs
${\approx}65$ (Chronos) to ${\approx}90$ (Toto) milliseconds of GPU time,
comfortably inside any 15--60-second control cadence, and the headline rests
on the 120M model. The real asymmetry is deployment footprint --- a
GPU-resident service versus arithmetic --- and I do not hide it; the closed
loop is what prices whether the margin earns it. The honest finding cuts
\emph{against} buying the biggest model: at provisioning cadence the point
accuracy edge over EWMA is only 2--9\%, and the SLA outcome is decided far
more by the quantile and the conversion than by which of the three TSFMs
supplies the band. What I will \emph{not} say is that the predictor is
outcome-irrelevant: the tested EWMA calibrations split on burst3 --- additive
and width-scaled fail, multiplicative passes only by carrying substantially
more standing work --- so the band's \emph{shape} is load-bearing even where
TSFM identity is not. That is precisely the value on offer: a calibrated band
a foundation model gives zero-shot without per-signal fitting --- exactly what
a multi-pool system, multiplying signals, needs. Even the cheapest of the
three suffices.

\Q{Your operational win is a calibrated p90 \emph{band} of a persistent signal, and your point forecast beats EWMA by only 2--9\%. You never ran the obvious cheap competitor --- a \emph{forward} EWMA read at its own empirical p90, through the same conversion and lead. How do you know the whole 0.62$\times$ is not obtainable for free, with the foundation model decorative?}
\A I ran exactly that cell --- a \emph{forward} EWMA at a p90 empirically
calibrated on the 1.4-day context window (three standard schemes: additive
conformal, multiplicative conformal, and a width-scaled native band, all at
eval coverage 0.88--0.90, matching the TSFM), through the \emph{identical}
issue-time pipeline --- and it does \emph{not} get the win for free. Under
the thesis's aggregate-and-per-episode grading the forward EWMA hits a
coverage-versus-cost bind: the two cheap calibrations fail the fast onset
(burst3 p99 1{,}285 and 2{,}084\,ms) and the high-volume onset (burst5,
backlog-divergent), passing only the \emph{burst-masked} aggregate
(0.70/0.56$\times$); the one calibration that clears every declared
evaluation episode
(multiplicative, burst3 \SI{747}{\milli\second}) costs \textbf{1.09$\times$
the static fleet} --- no economic win at all. Chronos clears every declared
evaluation episode (the finer daily-fold audit has the F2/F6 failures)
\emph{and} costs 0.62$\times$. So the forecaster is load-bearing: the TSFM's
zero-shot, state-adaptive p90 band is the one thing that covers fast onsets
\emph{without} calm-time over-provisioning, and a cheap forward forecaster's
calibrated band cannot be both at once.
\Point{This \emph{refines} ``model choice is second-order'': accuracy is
second-order (ARIMA ties Chronos), but the calibrated band \emph{shape} is
not --- and that shape is exactly what the TSFM supplies for free. Measured
(verification/results/ewma\_baseline\_*.json; chronos control reproduces the
0.62$\times$/\SI{589}{\milli\second} headline exactly), not asserted.}
\Concede{Scope it honestly: I tested three standard calibrations, not every
conceivable one --- a bespoke volatility-band EWMA might do better, so the
claim is ``standard forward-EWMA calibrations do not match the TSFM,'' not
``no cheap forecaster could.''}

\Q{Why not ensemble the three foundation models rather than pick one?}
\A Because ensembling optimizes the axis I showed does not matter. The three
are a noise-level tie on accuracy and pass the week at the \emph{identical}
\SI{651.1}{\milli\second} tail; the integer-replica conversion then
quantizes their differing forecasts onto the same worker schedule. An
ensemble would add inference cost and latency to shrink a spread the
conversion already erases, and the SLA outcome is set by the quantile and
the conversion, not the point forecast. If anything the finding argues the
other way --- deploy the smallest single model (120M Chronos), because model
choice is second-order.

\Q{SageServe already pairs a traffic forecast with a fleet optimizer and reports ${\sim}25\%$ GPU-hour savings. Strip the multimodal framing --- is your loop not just SageServe on a new signal?}
\A The signal \emph{is} the contribution, not incidental to it. SageServe
forecasts the aggregate \emph{request rate} of a text-only fleet and feeds
an ILP over multi-datacenter capacity, with a bespoke per-service
statistical predictor. In modality-disaggregated serving the request-rate
stream carries almost none of the vision pool's demand ($R^2{=}0.123$), so
``SageServe applied to the vision pool'' forecasts the wrong series and
reproduces the missizing I quantify. Two further separations: I use
pretrained \emph{zero-shot} models to kill the per-signal fitting burden a
multi-pool system multiplies, and I provision \emph{per pool} where
SageServe sizes one aggregate fleet.
\Trap{Do not litigate the numbers --- $0.62\times$ is encoder-only
worker-time, SageServe's 25\% is fleet GPU-hours; not comparable. Defend on
``which signal, and why the request stream is the wrong one here.''}

\Q{Dynamo's own planner --- the platform you simulate on --- grew a forecast-driven loop. Are you not re-implementing a feature of your own substrate?}
\A Dynamo's second, slower loop forecasts the aggregate request rate and
mean sequence lengths with classical models to set replica floors, and by
its own design notes it watches a single averaged text-side stream. I use
Dynamo's \emph{reactive} SLA loop unmodified as my baseline, so I am not
competing with its predictive loop --- I am showing what that loop cannot
do: forecast a non-rate, per-modality signal, with a zero-shot model rather
than a per-service fit. The overlap is a strength: my contribution could be
implemented \emph{as} an additional Dynamo forecast head, which makes it
deployable, not redundant.

% ==========================================================================
\section{Beat 3 --- CONTROL: placement and conversion, not clairvoyance}
\label{sec:beat3}

The mechanism beat, and the one whose story I most recently corrected. The claim is not that the forecast detects onsets. It is that forward placement of a calibrated band beats trailing placement at matched rating, and that turning demand into replicas against a measured tail knee rather than a mean is what decides pass from fail. Say \emph{forecast-anticipated} for vision and \emph{tail-aware conversion} for language, every time, including under pressure.

\Q{You claim the forecast ``detects every onset.'' Does it?}
\A No --- and that claim is gone from the manuscript. Under honest issue-time
actuation, only the clearly-gradual onset is genuinely \emph{anticipated}; the
faster ones are covered by \emph{graduated headroom} --- the p90 band standing
capacity up as the predictive distribution widens, not a crisp detection event.
The SLA result survives untouched (the \SI{588}{\milli\second} floor is
unaffected), but the mechanism is headroom, not clairvoyance. I corrected this
thesis-wide.
\Trap{This is the correction I am most likely to fumble, because the graduated-
headroom story is subtler than the detection story --- and I liked the detection
story. Rehearse it until the subtler version comes out first.}

\Q{Your central word is \emph{anticipatory}, yet you say the burst3 win comes from capacity standing ``several minutes ahead of the peak'' --- on a \emph{one-minute} onset, where you also say the forecast lags with no usable lead. Capacity that early can only be prior over-provisioning. On the one fast onset your story rests on, what does the forecast anticipate?}
\A On burst3, nothing about the onset --- I concede that cleanly, and it is
the correction I am most careful to make. The win is not foresight; it is
that the forecast's p90 band sits higher \emph{earlier}, during the ambient
activity before the spike, than a trailing-observed p90 does at matched
rating, so capacity is standing when the one-minute spike lands. Genuine
anticipatory lead is measured only on the gradual onset (burst4,
${+}\SI{345}{\second}$). The load-bearing comparison is therefore not
lead-over-ground-truth (the forecast has none on burst3) but
advantage-over-reactive-at-equal-$\rho$ (588 vs 1{,}973--11{,}568\,ms) ---
headroom \emph{placement}, not clairvoyance.
\Trap{Do not defend ``several minutes ahead of the peak'' on a one-minute
onset --- it is the sentence an examiner reads aloud, and it is a manuscript
wording I should fix. Say ``anticipation is demonstrated on burst4; burst3
is better-placed standing headroom.'' Note burst6 is also a one-minute onset
and \emph{is} reactively caught (its peak drains within a batch), so
burst3's difficulty is volume, not speed alone.}

\Q{Your only direct isolation of ``anticipation'' --- the reactive-with-headroom ablation --- uses a trailing window at matched $\rho{=}15$. On a one-minute onset no trailing window can place capacity in time \emph{by construction}, so the ablation is rigged to fail reactive. And if the burst3 win is ``the same over-provisioning the $\rho$-frontier prices,'' a sufficiently over-provisioned reactive floor should also pass. Where is that run?}
\A The ablation is a fair \emph{matched-rating} comparison, and I should name
what it matches: $\rho$, the signal, the conversion law, the hysteresis and
the lead --- not provisioned worker-time. On burst3 at $L{=}\SI{30}{\second}$
the forecast arm provisions 31{,}838.0 worker-seconds against the
trailing-p90 $W{=}480$ arm's 22{,}988.1, so equal $\rho$ is emphatically not
equal cost, and the reactive arm loses the onset while spending \emph{less}.
At that matched rating the forecast's projected p90 band beats a trailing p90
band, and headroom alone (which cuts the reactive tail up to $8\times$) still
fails, so headroom is necessary-not-sufficient \emph{at matched rating}. That
is the honest claim: at equal $\rho$ and mechanics, forward placement beats
trailing placement on fast onsets.
\Concede{``At equal $\rho$'' is not ``anticipation is necessary,'' and the
brute-floor run exists: the $M{=}1\ldots32$ family was swept. Under the
post-hoc common-warm diagnostic $M{=}2$ passes every evaluation window at
$0.746\times$ static ($1.20\times$ the forecast) and the calibration-frozen
$M{=}4$ passes at $0.887\times$; under the shipped strict re-warm mechanics
the cheapest passing floor is $M{=}6$ at $1.042\times$ static. So concede the
demotion outright rather than offering the run: the claim is ``anticipation
makes the same SLA cheaper,'' not ``anticipation enables the SLA'' --- still
true, and weaker, and the floor sweep is reactive-favorable post hoc while
the forecast stayed one frozen point.}

\Q{You built the encoder reactive baseline yourself. How do I know you did not hobble it?}
\A Because it is \emph{identically informed} with the forecast policy, and I
separate the remaining dials by ablation. It gets the same worker-mapping
law, the same $\rho{=}15$, the same 3-minute hysteresis, the same cold-start
lead, and the correct per-pool signal --- the observed image rate, not the
request rate --- read reactively rather than forecast. The only differences
left are the p90 headroom band and where that headroom is placed ---
forward-projected versus trailing --- and I isolate them: a
reactive controller \emph{given} the tail-aware p90 headroom of the observed
rate still fails the one-minute onset at every lead (p99
$1{,}973$--$11{,}568$\,ms against the forecast's \SI{588}{\milli\second}),
even though that headroom cuts its tail up to $8\times$. So at matched
\emph{rating and mechanics} the baseline is not headroom-limited
(reactive-with-headroom fails) and not
signal-limited (it reads the right signal); the residual is forward
\emph{placement}. And I will state the limit of that matching before anyone
asks: matched $\rho$ is not matched cost. On burst3 at $L{=}\SI{30}{\second}$
the forecast arm provisions 31{,}838.0 worker-seconds against the trailing-p90
$W{=}480$ reactive arm's 22{,}988.1 --- the reactive arm is the
\emph{cheaper} one there and still loses the onset, which is why I say
matched rating rather than equal cost. Spend does buy the tail, and I say so: with a deep enough
standing floor the reactive family passes --- post-hoc common-warm $M{=}2$ at
$0.746\times$ static, calibration-frozen $M{=}4$ at $0.887\times$ --- just
never as cheaply as the fixed forecast point. The language-side reactive baseline
is stronger still --- Dynamo's production SLA planner in both shipped
variants, unmodified.
\Point{The reactive-with-headroom ablation is the whole answer to
``strawman.'' Lead with it.}

\Q{Why is mean-based capacity conversion ``the defect''? How do you know it is not just bad tuning?}
\A Because I ran the triad that separates them. On the same burst window
(burst1) at a realistic 60-second replica lead: reactive fails
(\SI{1616}{\milli\second}); the forecast
fed through the \emph{mean} conversion also fails, at essentially the same tail
(\SI{1717}{\milli\second}); the forecast with a tail-aware conversion passes
(\SI{736}{\milli\second}). Better signal through the wrong conversion buys
nothing. And tightening the reactive controller's targets under lead raises spend
from 18.6 to \SI{20.4}{GPU-h} while the tail stays pinned around
\SI{1620}{\milli\second} --- the budget is not even saturated. So it is not
thrift and it is not tuning; it is information timing and conversion.
\Point{The triad is the cleanest single piece of evidence in the thesis for the
central mechanistic claim. Lead with it whenever the mechanism is questioned.}

\Q{What is the language-pool result, exactly?}
\A This one I want to state precisely, because an early draft over-claimed it.
The vision pool is \emph{forecast-anticipated}: its signal is not rate-visible,
and the forecast stands capacity up before the peak. The language pool is
\emph{not} a second anticipation experiment. Its signal \emph{is} rate-visible,
and what the experiment tests is the \emph{conversion} --- turning demand into
replicas against a measured tail knee rather than a mean --- under a held cursor.
Both ingredients are necessary, but they are different claims, and I narrowed the
manuscript thesis-wide to say so.
\Trap{An early draft said ``each pool on a forecast of its own signal.'' That
over-claims the language side, and an examiner who read the code will catch it.
Say \emph{forecast-anticipated} for vision, \emph{tail-aware conversion} for
language --- every time, including under pressure.}

\Q{Your tail-aware conversion multiplies the planner's replica belief by ${\approx}6\times$ ($\hat r/r_{\text{tail}}=10.07/1.62$), calibrated on two windows by single-replica flips. How robust is that?}
\A The multiplier is not a tuned constant --- it is the ratio of the
planner's throughput belief to the \emph{measured} tail knee
$r_{\text{tail}}$, the largest offered rate a replica sustains under the
tail SLA (1.62\,req/s on burst1, bracketed $[1.26,1.44]$ on the binding
burst4 mix). Validation flips exactly at the formula's boundary on both
windows tested (burst1: 4.64 fails, 6.22 passes; burst4: 6.22 fails, 7.0
passes), and a ${\approx}14\%$ margin above the binding knee is policy
precisely because $r_{\text{tail}}$ is mix-dependent and planner-in-loop
replays show ${\approx}4\%$ p99 run noise.
\Concede{The $6\times$ gap is large and the end-to-end calibration is two
mixes at single-replica granularity; the remaining windows are priced by the
knee formula, not replayed. Defend ``measured knee, not tuned constant''
plus the 14\% margin; own the two-window scope.}

\Q{Why p90 as the provisioning quantile? Is it newsvendor-optimal, or just a round number?}
\A A critical-fractile choice, and I treat it as one. Which quantile to hold
is a newsvendor problem (Arrow, Harris \& Marschak 1951): the optimal level
follows from the ratio of under-provisioning cost (SLA damage during a
cold-start lead) to over-provisioning cost (idle GPU-seconds). Rather than
fix p90 by convention I sweep p50 against p90 under the SLA and let the
outcome choose --- p90 passes where p50 fails on fast onsets, because a
median forecast sizes to the mean and floods the FIFO. And p90 is not a
generic margin: its upper edge is empirically \emph{calibrated} at
provisioning cadence (coverage 0.88--0.90; 11\% exceedance against a nominal
10\% on the largest burst), so ``provision at p90'' concedes bounded,
calibrated exceedances, not an arbitrary cushion.
\Point{Newsvendor is the principled answer to ``why p90.'' Pair it with the
calibration: p90 is a \emph{measured} edge, not a guess.}

\Q{Your p90-headroom rule looks like square-root staffing. Why not just use Erlang machinery?}
\A It is the same instinct --- provision against demand uncertainty, not
demand mean --- and I cite that lineage deliberately (Halfin--Whitt 1981;
Borst, Mandelbaum \& Reiman 2004). But classical staffing needs what this
workload refuses to provide: (near-)Poisson arrivals and a stationary rate,
so that $\sqrt{\lambda}$ safety capacity covers the fluctuations. Here the
``fluctuation'' is a $16\times$-mean flood arriving under long-range
dependence (Fano 938 at the hour), where square-root headroom on the trailing
mean is off by an order of magnitude exactly when it matters. The forecast's
p90 band is the workload-adaptive generalization: it \emph{is} the
uncertainty term, estimated from the signal's own conditional distribution
rather than from a Poisson assumption, and it is empirically calibrated
(coverage 0.88--0.90 at nominal 0.90). Where classical theory does apply I
use it --- the service side is a plain FIFO with a measured batch service
curve --- the departure is confined to the arrival side, where the data
forces it.

\Q{The burst breaks the encoder pool and you say ``only capacity fixes it.'' Would a smarter queue discipline --- shortest-job-first, priority for small requests --- not remove the blocking without buying GPUs?}
\A Scheduling and provisioning are orthogonal levers, and this thesis is
deliberately about the second. A reordering discipline reshapes \emph{which}
requests wait during a flood, but the flood's total work is fixed: the
worst second lands 2{,}726 images and a frames episode enqueues tens of
thousands, so a p99 over all requests still lands on frames work (2.7\% of
requests, above the 1\% tail). Chunk-parallel dispatch already spreads an
oversized request across free workers; under the sequential alternative the
floor is capacity-bound at any discipline. And I say explicitly that
minute-level forecasts are \emph{provisioning} inputs, not scheduler
signals --- dispatch policy and capacity handle the second-scale blocking,
the forecast decides how much capacity is standing when the flood lands.
\Concede{I did not sweep alternative disciplines (SJF, priority, PS); FIFO
is the modeled discipline. The claim is that aggregate flood work needs
capacity, which reordering cannot conjure.}

\Q{Why these SLA thresholds --- TTFT and encoder-delay p99 under \SI{1}{\second}, ITL p99 under \SI{100}{\milli\second}? Are they not arbitrary?}
\A A design choice, which I both justify and stress-test. The values are
exactly twice Dynamo's own default planner targets (\SI{500}{}/\SI{50}{\milli\second},
which the planner applies as window-\emph{means}); I apply them as
\emph{tail} bounds --- first token within a second, at least ten
tokens/second for the worst-percentile user --- an interactivity envelope
consistent with disaggregated-serving evaluations (DistServe, Chiron). The
sweep curves themselves are SLA-free and re-gradable under any bounds, and
I report where the verdict moves: under a lax TTFT-only \SI{2}{\second}
bound the language knees sit one to two octaves higher, but the vision-pool
queueing verdict and the TP1 prefill floor survive threshold changes.
\Trap{Do not defend \SI{1}{\second} as canonical. Defend it as $2\times$ the
platform's own defaults, applied as tails, with the sweep re-gradable ---
and note the vision verdict is threshold-robust.}

\Q{You grade encoder-delay p99 under \SI{1}{\second} and language TTFT p99 under \SI{1}{\second} separately. But an image request waits for both. What is the end-to-end first-token SLO?}
\A There is no combined budget, and I will say so first. The thesis grades
each pool against its own bound and never sums across hardware classes,
because the pools are provisioned by different controllers on different
silicon and per-pool attribution is the object of the study (either pool
replays in isolation). The consequence is honest and stateable: since the
encoder-delay p99 already includes the \SI{272}{\milli\second} handoff and
the language TTFT p99 is a \emph{separate} 651--\SI{736}{\milli\second}, a
serial image request can see roughly \SI{1.2}{\second} to first token at the
joint tail even when both pools ``pass.'' The defensible position is that a
per-pool objective is the \emph{actionable} one for a per-pool provisioner
--- but I claim no bounded end-to-end experience.
\Trap{Do not present the \SI{588}{\milli\second} floor as the served
first-token latency --- it is encoder-side only. Concede the end-to-end
budget is undefined and frame per-pool as the provisioning objective.}

\Q{Your SLA sets two separate 1-second bounds, and you meet both --- yet your own reconstruction puts ingress-to-first-token p99 at ${\approx}\SI{1.30}{\second}$, 30\% above every ``1\,s'' a reader sees. ``Holds the 1\,s tail'' is per-pool bookkeeping; the user waits 1.3\,s. Where do you say so?}
\A I say it once, and I concede it should be promoted: the two per-stage p99s
do not compose to \SI{1}{\second} because the encoder and prefill tails are
positively correlated ($r{=}0.52$ --- the same frames requests drive both),
so the union lands near the sum (${\approx}1.24$--\SI{1.30}{\second}), not
the max, and ``the two separately reported percentiles understate the
modelled first-token wait.'' It sits inside the ${\approx}\SI{2}{\second}$
union budget I state (a union bound at ${\approx}$p98, not an end-to-end p99
guarantee), and the serialized seam is conservative (overlap could only
lower it).
\Concede{Promote the ${\approx}\SI{1.30}{\second}$ ingress-to-first-token p99
into the abstract/conclusion, framed as ``$\sim2\times$ either marginal from
$r{=}0.52$ tail coupling, within the 2\,s union.'' Do not defend two 1-second
bounds as if they imply one second to first token --- they provably do not.}

\Q{Your replay is open-loop --- arrivals ignore experienced latency. Under real overload clients abandon and retry, draining the very queue that sinks the reactive baseline. Does open-loop not inflate the reactive failure?}
\A It might, in either direction, and I flag exactly this. The defense is
that offered load is applied \emph{identically} to every policy, so the
comparison is fair even if the absolute overload is pessimistic; the claim
is comparative under fixed offered load, not a prediction of realized
behavior under admission control. I concede the direction the question
pushes --- abandonment preferentially relieves the overloaded (reactive)
policy, so a closed-loop model could narrow the gap, though retries could
equally widen it. The absolute magnitude of the reactive failure (burst3,
9.5--\SI{54}{\second}) is an upper bound, not a realized user experience;
the \emph{direction} of the verdict is what I claim.

\Q{Two of the six bursts are excluded from forecast grading. Are you not scoring on the windows that flatter the forecast?}
\A The exclusion is a zero-shot honesty constraint, not a selection. The
first 1.4 days are consumed as the models' rolling history before their
first usable forecast; burst2 and burst6 fall inside it, so they are graded
by every \emph{forecast-free} baseline (static grid, reactive) but not by a
forecast policy, because the forecast is structurally absent there --- a
hole a deployed system entering with prior history would not have. The
fast-onset stress case is still in the forecast evaluation: burst3, the
one-minute onset, is in-span, is the hardest cell, and is the one carried
live.
\Concede{The $0.62\times$ is scored over four in-span bursts, not six, and
burst6 --- a fast onset the reactive baseline actually catches --- is never
put to the forecast policy. Say ``four in-span, burst3 carries the fast
onset,'' not ``all six.''}

\Q{What would falsify your central claim?}
\A Cleanly: a production multimodal trace in which image demand \emph{is}
recoverable from request rate --- an $R^2$ near 1 rather than 0.12 --- would kill
the premise outright. Short of that: if the image signal's persistence turned out
to be an artifact of this trace, and the signal proved unforecastable elsewhere,
the escape route closes and reactive scaling is all there is. And on the mechanism
specifically: if a reactive controller watching the \emph{correct} signal could be
tuned to hold the tail through the fast onsets \emph{as cheaply}, then the
forecast machinery would be unnecessary and the contribution collapses to
``watch the right metric.'' I tested that directly --- the
reactive-on-the-same-signal ablation at matched rating --- and it
fails at every lead. And the stronger variant too: a reactive controller
granted the \emph{same} tail-aware p90 headroom it was supposedly denied
still fails the one-minute onset at every lead and window (p99
$1{,}973$--$11{,}568$\,ms), even though that headroom cuts its tail by up to
$8\times$. What I will \emph{not} claim is that reactive can never pass: the
whole $M{=}1\ldots32$ warm-floor family was run, and under the post-hoc
common-warm readiness diagnostic $M{=}2$ passes every evaluation window at
$0.746\times$ static while the calibration-frozen $M{=}4$ passes at
$0.887\times$. Both cost more than the fixed forecast point. So the
established mechanism is cheaper \emph{placement} of standing headroom at
matched rating, not the necessity of forecasting for feasibility --- and the
falsifier is correspondingly ``a reactive point that holds every evaluation
window at no more than the fixed forecast point's worker-time.'' That is a
distinct experiment from the untested symmetric forecast-plus-floor family,
which I also owe: the first would falsify the observed tested-point gap, the
second only repairs design-space symmetry.
\Trap{This question rewards a specific falsifiable answer and punishes a vague
one. Do not say ``more traces would help.'' Name the observation that kills the
claim.}

\Q{You cite BlitzScale, which multicasts weights to make scale-up near-live. If cold starts are near-eliminated, does a reactive loop not suffice --- killing your motivation?}
\A BlitzScale attacks \emph{weight loading}, which is the language pool's
lead, not the vision pool's. The 576M encoder's checkpoint is
\SI{0.11}{\second} to read and \SI{0.13}{\second} to transfer; its
${\approx}3$--\SI{5}{\second} of model-side work is interpreter and
framework import, and encoder-pool lead above ${\approx}\SI{5}{\second}$ is
\emph{orchestration}-bound (pod scheduling, image pull, node acquisition) ---
which weight multicast does not touch. Against a one-minute onset with a
\SI{208}{img/s} single-minute jump, even a near-live weight path leaves the
orchestration lead exposed, which is why the closed loop sweeps leads from
15 to \SI{165}{\second} and reports lead-invariance across the realizable
range rather than tying the result to one pessimistic constant.
\Trap{Do not argue BlitzScale is irrelevant; argue it moves the
\emph{language}-pool lead, and the burst breaks the \emph{encoder} pool.}

% ==========================================================================
\section{Beat 4 --- ECONOMICS: the same compliance, less standing capacity}
\label{sec:beat4}

The payoff beat, and the one with the most audit history, so it is the beat to answer precisely and leave. The claim is a trade-off with an identified failure mode, not universal superiority: equal declared compliance at less provisioned worker-time, encoder-side only, in two per-class ratios that are never summed.

\Q{Walk me through where the 0.62x comes from.}
\A Numerator first, denominator second, and one word --- \emph{issue-time} ---
doing the honesty work. The numerator is the forecast-driven encoder
schedule's provisioned worker-seconds over the covered 5.60-day span:
2{,}410{,}755, generated by the p90 quantile path at rating $\rho{=}15$ with a
30-second cold-start lead, actuated on forecast \emph{issue} time --- each
scaling decision uses only forecasts that had actually been issued by the
moment it acts, never a forecast dated to its target. The denominator is the
cheapest \emph{deployable} static fleet: $C{=}8$ encoder workers held all
week, 3{,}869{,}278 worker-seconds, itself verified week-feasible (p99
\SI{644}{\milli\second}, backlog slope ${\approx}0$). Both schedules replay
through the identical converter over the identical span; both meet the SLA;
the ratio is $2{,}410{,}755 / 3{,}869{,}278 = 0.62$. It is nearly
lead-invariant across the archive's realizable issued leads (2.41M
worker-seconds at both 15\,s and 165\,s) under the warm-pool readiness model,
and 0.59--0.62$\times$ under a strict re-warm model --- both disclosed.
\Point{If asked ``why is the denominator $C{=}8$ and not something bigger'':
because bigger was my \emph{error} --- the original $C{=}16$ denominator came
from an endpoint-divergence artifact in one window, and correcting it is
exactly what moved the headline from $0.31\times$ to $0.62\times$ against my
own interest.}

\Q{The headline was once 0.27--0.31x. Why did it move to 0.62x?}
\A Two corrections, both found through adversarial review, and both of which I
would rather state than have extracted. First, the denominator was wrong: the
static baseline had been sized at $C{=}16$ because of an endpoint-divergence
artifact in one burst window; the honest deployable static fleet is $C{=}8$,
which halves the denominator. Second, the forecast was being actuated on
\emph{target} time rather than \emph{issue} time --- the schedule was implicitly
allowed to use a forecast before it would actually have been available. Rebuilt
on honest issue-time actuation, the number is $0.62\times$. What survived both
corrections is what matters: the \SI{588}{\milli\second} service floor, every SLA
verdict, and every policy-versus-policy direction. The bug cost me the size of
the win, not the win.
\Trap{This is the single best opportunity in the defense to show that I audit my
own work harder than the committee will. Say the number moved, say why, say what
survived --- in that order, without defensiveness.}

\Q{Is rho=15 cherry-picked to make the number look good?}
\A No, and I can show the frontier rather than assert it. The result is a
tunable frontier, not a point: p90 passes at every rating from $\rho{=}10$ to
$\rho{=}20$, at costs from $0.89\times$ down to $0.48\times$. I report
$0.62\times$ at $\rho{=}15$ as the \emph{conservative} operating point --- it
holds the \SI{588}{\milli\second} floor with margin. The cheapest clean rating is
actually $\rho{=}20$ at ${\approx}0.48\times$, which I disclose; I did not quote
it because its per-window p99 climbs toward the bound under the strict readiness
model. If anything I am reporting a number worse than the best one available.
\Concede{$\rho{=}15$ was chosen \emph{ex post} as conservative. I say so. What I
claim is not that nothing was tuned, but that the full frontier is reported so
the choice is auditable.}

\Q{You report 0.62$\times$ on the encoder pool and 0.965$\times$ on the language pool and never add them. The encoder runs cheap consumer 5090s; the language pool runs H100s that dominate real spend. Priced end-to-end in one currency, is your headline not a small cut on the cheap line while the expensive line is break-even --- or 1.16--1.19$\times$ with Toto or TimesFM?}
\A The two pools carry two different claims, and the language one was never
cost. The encoder claim is cost (2.41M vs 3.87M worker-seconds, 0.62$\times$).
The language claim is \emph{correctness at parity cost}: the reactive planner
fails the week (TTFT p99 \SI{1197}{\milli\second}) while spending \emph{less}
(281.8 vs 518.6 GPU-h) --- ``spend is not the defect.'' Pricing a compliance
result in dollars is a category error, and blending 5090-hours with
H100-hours needs a price ratio the trace cannot supply, which is why the
units stay separate and the thesis statement says so.
\Concede{There is no end-to-end TCO figure, and ``language near-break-even''
is Chronos-specific --- with Toto/TimesFM it is 1.16--1.19$\times$, a net
loss. Concede the economic headline is scoped to the encoder pool's
worker-time, and state ``0.965$\times$ (Chronos); 1.16--1.19$\times$ under
other forecasters'' rather than let it be extracted.}

\Q{Your deployable static fleet --- 8 cheap 5090s --- already \emph{passes} the SLA at 1.00$\times$. So the naive operator buys 8 cards and never thinks again. Your machinery saves ${\approx}3$ average consumer cards. Is a foundation-model pipeline solving a problem anyone would pay to solve?}
\A ``Just buy 8'' quietly assumes you know the peak --- and you do not:
minute-resolution monitoring recovers under 4\% of the true instantaneous
image peak, which is ${\approx}16\times$ the mean and rate-invisible. The
static $C{=}8$ baseline is an \emph{oracle}, sized from the whole week's
observed maximum, knowledge a live operator lacks a priori; the forecast
reaches 0.62$\times$ of that oracle under \emph{causal actuation} --- every
schedule is issued from a rolling-origin forecast that never reads its own
future. What is not causal is the choice of operating point: the headline
rating $\rho{=}15$ was selected \emph{ex post}, with the whole frontier
disclosed, so ``only causal information'' is a statement about runtime
actuation, not about how the arm was picked. And
over-provisioning fails \emph{invisibly}: the pool violates its bound at
29\% mean utilization via episodic blocking a dashboard cannot see, so ``set
it and forget it'' is unsafe, not just wasteful --- you cannot pick the right
constant without characterizing the peak forward in time, which is the
forecasting problem.
\Point{Reframe 0.62$\times$ as ``near-oracle sizing from causal
information,'' not ``38\% off a number you already had.'' The choice is
forecaster-vs-guess, and the guess fails from an incident you cannot see
coming.}
\Concede{For one small cluster the absolute dollars are modest; the
efficiency case scales with fleet size and encoder growth, which I do not
quantify beyond one trace.}

\Q{Could the operator not just buy 2x capacity and skip the machinery?}
\A They can, and the static $C{=}8$ baseline \emph{is} that policy priced
honestly --- it passes the week. The thesis's claim is not that static
provisioning fails; it is that it costs $1.6\times$ the forecast policy for
the same SLA. The reactive alternative that promises static's savings without
forecasting does exist, but only with standing capacity bought back: at
matched rating and no floor it fails the fast onset at every lead, and the
floors that do pass cost more than the forecast point --- post-hoc
common-warm $M{=}2$ at $0.746\times$ static ($1.20\times$ the forecast),
calibration-frozen $M{=}4$ at $0.887\times$, and under the shipped strict
re-warm mechanics the cheapest passing floor is $M{=}6$ at $1.042\times$
static. So the real menu is: pay $1.6\times$
for simplicity, roughly $1.2\times$ for a floor whose depth you can only pick
after the fact, or run the loop. At encoder-pool scale that gap is real money
(the static fleet holds 8 workers around the clock against the policy's
demand-following 2.41M worker-seconds), and the machinery is one 120M-model
forecast pass per control tick. The economics get more compelling, not less,
as the pool grows.

\Q{You price the encoder pool on RTX~5090 and the language pool on H100 in one system. Is a cross-pool cost comparison not apples-to-oranges?}
\A The mixing is the architecture, not an accident: modality disaggregation
exists precisely so each pool runs on its own hardware class, and I never
sum GPU-time across classes --- the headline is two separate per-class
ratios ($0.62\times$ encoder, $0.965\times$ language). The
``order-of-magnitude apart'' knee result is not a raw GPU-time comparison
either: it grades each pool's saturation against the \emph{same}
\SI{1}{\second} tail budget on the same offered-load axis, and the
separation reflects demand structure (few, cheap, well-amortized images per
request versus long decode-bound prefills), not the hardware classes --- at
the language knee the encoder idles at 4.5\% utilization. Putting both on
one GPU would defeat the premise under study.

\Q{With two of your three forecasters the language pool costs \emph{more} than the static fleet (Toto $1.16\times$, TimesFM $1.19\times$). Why is the language pool in the economic story at all?}
\A It essentially is not, and I present it as a mechanism result, not a
savings one. The decisive per-pool economics are on the encoder side; the
language pool's dynamic saving is structurally small because replica scaling
cannot morph a burst-required TP2 shape into a calm-day TP1 one, and the
$0.965\times$ is explicitly Chronos-specific (the pass/fail \emph{verdict} is
model-independent, the cost is not). What the language pool establishes is
that the \emph{conversion} generalizes across pools: fed the planner's
native mean-based conversion, the same forecast \emph{fails} the pool
entirely (the 736 vs 1{,}717/1{,}616\,ms triad). Its only economic
contribution is cross-window (TP1$\times$2 on calm traffic).
\Point{Keep the economic headline strictly encoder-side. If pressed for an
all-in system number, there is none --- two per-class ratios, never summed,
because a 5090-second and an H100-second are not fungible.}

\Q{ModServe already reports 25--41.3\% serving-cost savings from modality-aware scaling. Your headline is 0.62$\times$ --- 38\%, inside their band. Did you re-derive their number?}
\A Different baselines, so the coincidence is arithmetic, not substance.
ModServe's 25--41.3\% is the gain of modality \emph{disaggregation} itself
--- fewer GPUs at SLO versus a monolithic/request-rate baseline. My
0.62$\times$ is the gain of \emph{forecasting} on top of an
already-disaggregated, modality-aware \emph{static-optimal} fleet --- a
strictly harder baseline. They stack; they do not compete. And I claim no
architectural novelty: the systems that ship the encoder-scaling knob all
scale it \emph{reactively} on a stream carrying $R^2{=}0.123$ of the
encoder's demand, and ElasticMM's own authors call the signal ``difficult
to predict.'' My contribution is that it is not.

\Q{The obligatory cost/energy/sustainability question.}
\A The thesis is, at bottom, a waste-reduction result, and I can state it in
those terms: meeting the same tail SLA at $0.62\times$ the provisioned
worker-time is a ${\approx}38\%$ cut in encoder fleet-time, and the
$\rho$-frontier shows an operator willing to run nearer the bound can reach
$0.48\times$. The mechanism generalizes: static fleets are sized for peaks
that occupy a small fraction of the week (the burst episodes cover minutes to
hours of a 5.60-day span), so any anticipation-based policy converts tail
insurance from a standing cost into an on-demand one. Two honest caveats:
GPU-time is a proxy --- I do not measure watts, and idle-versus-active power
means worker-seconds overstate the energy delta; and the forecasters
themselves cost compute, though the headline rests on a 120M-parameter model
whose inference is noise against a single H100-hour.

% ==========================================================================
\section{Beat 5 --- VALIDITY: what is measured, what is modeled, and where each is anchored}
\label{sec:beat5}

The fidelity beat. Every answer here begins by saying which side of the measured/modeled seam a number came from. The decomposition is deliberate --- silicon for the encoder cost, a live control plane for the planner, a sensitivity bracket for the H100 timing --- and the honest concession is that no single experiment puts real compute, the real control plane, and real weights on the request path at once.

\Q{What is the weakest part of the work?}
\A The economics --- specifically, that the week-scale cost result is
\emph{simulated}. The descriptive core is measured (a real production trace), the
cost model is measured (real silicon, matched engine, dense sweep), and the
control-plane verdict is replayed through the real Dynamo control plane --- but
with mocked engines, on one burst window, with a single run per policy. The
$0.62\times$ over 5.60 days is a simulation result, priced with a
silicon-validated cost model and a queue-replay under measured cost, but nobody
ran a GPU fleet for a week. I would not want the committee to hear that from
them first.
\Trap{Do not try to make this sound smaller than it is. Own it plainly, then say
precisely what \emph{is} measured --- which is a great deal.}

\Q{How do I know the simulator is not just confirming its own assumptions?}
\A Three independent anchors, and I would rather walk through all three than
assert the point. First, the encoder cost model is not a fit --- it is
performance data I collected on the silicon itself, at locked clocks, inside the
pinned production container, and upstreamed. Second, the dominant operator agrees
with physical measurement within ${\approx}3\%$ and the end-to-end kernel ledger
closes with no unexplained residual. Third, and most important, the composed cost
model was validated against \emph{the serving engine's own vision tower} (vLLM)
over a dense sweep of 69 batch sizes, holding within $12.8\%$ across the entire
$N \le 64$ regime the pool actually batches into. And the closed-loop result was
re-run with the measured vLLM cost substituted for the model's: the p99 moved
from \SI{589}{\milli\second} to \SI{606}{\milli\second} and the provisioned
worker-seconds did not move at all.
\Point{The queue-replay control/treatment pair: \texttt{VLLM\_RATIO=0} reproduces
the committed \SI{588.37}{\milli\second} baseline at \SI{589.30}{\milli\second}; \texttt{VLLM\_RATIO=1} gives
\SI{605.61}{\milli\second} at the identical 2,410,754.7 worker-seconds. The cost
is the \emph{schedule's}, not the cost model's.}

\Q{Why simulation rather than a real cluster?}
\A Because the experiments require things a production cluster cannot give.
They sweep offered load \emph{past saturation}, replay the same week under many
configurations and policies, and --- decisively --- require counterfactuals: the
same trace, the same arrivals, a different scaler. You cannot run a
counterfactual on a live cluster; the arrivals are gone. Trace-driven
discrete-event simulation is the established instrument at this scale in LLM
serving. What I did \emph{not} do is treat fidelity as an assumption: every
timing model is either collected on silicon and physically validated, or
compared against its alternative so its effect on the conclusions is known.

\Q{The live validation --- how strong is it really?}
\A Weaker than the word ``validation'' suggests, and I have narrowed the claims
accordingly. What was carried live: the p90-forecast policy and the reactive
baseline, on the fastest-onset burst, through Dynamo's real control plane with
\emph{mocked engines}. Both verdicts reproduce in both directions --- forecast
passes live (695 versus \SI{651}{\milli\second} offline), reactive fails live
(3531 versus \SI{3495}{\milli\second}) --- with TTFT within $1.1$--$6.8\%$ and
replica-time within $8\%$. What was \emph{not} carried live: the p50 leg (offline
only) and any repeated trials. Those are point differences, not a noise band. I
say ``directionally, within $1.1$--$6.8\%$,'' never ``within run noise,'' because
the latter implies a variance estimate I do not have.
\Concede{One run per policy, no repetitions. Say it without being asked.}

\Q{Your ``live'' validation runs mocker workers loading the AIC latency model --- no GPU, no real forward. What does it validate that the offline replay did not?}
\A The control plane, and I scope it to exactly that. What becomes real that
was virtual offline: the HTTP frontend, the in-process KV router,
independent mocker worker \emph{OS processes} discovered over etcd (so
replicas die independently), Prometheus scraping, the SLA planner as its own
service, and actuation through the VirtualConnector protocol via a
worker-manager written for this thesis. Engine \emph{timing} is identical by
construction, so this is not a silicon check --- that is the separate
RTX~5090 anchor. The payoff is concrete: the live stack surfaced three
control-plane divergences invisible offline (inverted \texttt{startup\_time}
semantics, a capacity-model starvation at calm load, tick-cadence stall),
all harness properties, not policy ones.
\Concede{Neither anchor validates compute and control simultaneously; two
experiments cover two layers, and that combined test is named future work.
Do not oversell ``live.''}

\Q{Your one-breath says ``validated on a live control plane.'' That plane covers only the language pool --- and the language result is \emph{not} anticipation (a held-cursor conversion test), is break-even-to-loss (0.965--1.19$\times$), and its burst pass is token-convention-dependent. So the only pool you carried live is the weakest, and the vision pool the thesis rests on is never live-validated. Is your live anchor attached to everything except the contribution?}
\A The split is by necessity: the language pool is the only pool with a
\emph{production} control path (Dynamo's SLA planner), so it is the only one
that \emph{can} be live-validated on this stack; the vision encoder pool has
no shipped controller, so its anchor is necessarily the silicon cost model
--- a different, independently strong anchor. The live run validates
control-loop \emph{fidelity} (offline$\leftrightarrow$online within 1.1--6.8\%),
evidence about the platform, not a claim that the language pool is the
contribution.
\Concede{The live anchor validates the \emph{conversion}, not anticipation,
and its economics are Chronos-specific. The contribution --- anticipatory
vision provisioning --- is anchored on silicon cost, not a live loop, and a
live encoder loop is the named Phase-E gap. Do not let ``validated live''
stand unqualified; the follow-up is ``validated which claim?'' and the honest
answer is ``not the central one.''}

\Q{Point me to the single experiment where real GPU compute, the real control plane, and real model weights are all on the request path at once. The encoder anchor is real silicon but a standalone forward on \emph{random weights}; the live anchor is the real control plane but \emph{mocked engines}; the language timing is a vendor database. Three corners, empty middle --- why believe the composition?}
\A The decomposition is deliberate: validate each layer where it is cheapest
and where simulation cannot reach --- silicon for the encoder cost (the layer
AIC could not price at all; its SOL fallback is $5.9\times$ optimistic), a
live control loop for the planner (both verdicts reproduce, TTFT within
1.1--6.8\%), a sensitivity bracket for the H100 timing. The ``shape-determined,
so random weights suffice'' assumption is sound for the two dense operators I
collect (non-causal FlashAttention, bf16 GEMM: no data-dependent sparsity),
and the kernel ledger closes to 3.2\% on attention --- but I do not claim it
for the full forward under real weights and sustained-serving thermals.
\Concede{No single all-real experiment exists, and the specific missing rung
is a weight-loaded vLLM encoder worker in the queue --- named as future work.
Also concede $n{=}1$ die: I cannot error-bar another 5090, a contended run
swung the ratio $2\times$, and the clock-invariant \emph{ratio} protects
against boost differences, not bin or thermal ones.}

\Q{The abstract sells a ``calibrated H100 language simulator'' and RQ3 answers ``platform validity: yes, with measured fidelity.'' But Chapter~5 says the language timing ``cannot be re-validated against physical H100 within this work,'' and the live mockers load the \emph{same} AIC model. So no physical measurement backs any language latency that decides a verdict --- and even ``validated physically'' on the vision side spans 3\% (one operator), 12.8\% (the composed forward), and 0 for language.}
\A The asymmetry is the right one and disclosed in RQ3: the encoder cost model
is my own contribution with no prior data, so it was validated on silicon;
the H100 language timing is NVIDIA's shipped, widely-used database, taken as
given like any vendor spec, with the polynomial-vs-AIC delta bounding the
\emph{conclusions'} sensitivity to it (the burst verdict flips, the steady
knee moves an octave). I should lead the vision claim with the model-level
12.8\%, not the operator-level 3\%.
\Concede{``Calibrated'' is carrying the reader past ``unvalidated vendor
timing,'' and the frames-class prefills that decide the burst are priced by
\emph{extrapolation} past the grid edge --- so the language verdict is
convention- \emph{and} extrapolation-dependent. Rewrite ``validated
physically'' to ``encoder cost validated on silicon (12.8\% composed
forward); language timing from AIC's shipped H100 database.''}

\Q{The encoder cost model is silicon-validated, but the language pool's H100 timing is not. Is half the two-pool result resting on an unvalidated model?}
\A Yes, the fidelity is asymmetric, and I put that in print rather than
hide it: the language pool's timing is trusted from AIC's shipped
\texttt{h100\_sxm} database and Dynamo's engine and, within this work,
cannot be re-validated against a physical H100. Three things bound the
exposure. First, I claim \emph{sensitivity}, not accuracy, for that side
--- the calibrated-versus-polynomial comparison quantifies how far the
conclusions move under the timing model, and its \emph{direction} (the
default polynomial flattens the tails, missing a $1.9\times$ burst-window
violation outright) is robust even where its magnitude is not. Second, the
load-bearing language verdict is a compute-bound \emph{floor}: a
${\approx}36$k-token TP1 prefill busts the \SI{1}{\second} TTFT budget at
\emph{any} worker count, which is arithmetic, not a calibration subtlety.
Third, the thesis's central contribution is the \emph{vision} result, which
rests on this thesis's own validated silicon and is independent of the H100
model.
\Concede{The language side is a calibrated-model sensitivity study, not a
measured one. Say it before it is said to me.}

\Q{Is a 12.8\% cost-model error acceptable? That sounds large.}
\A The disciplined answer has three steps, in order. First, the error is
\emph{bounded and measured}, not estimated: a dense sweep of 69 batch sizes,
$N{=}1$ through 512, against the serving engine's own vision tower, with the
maximum deviation $+12.7\%$ at $N{=}52$ inside the $N \le 64$ regime the pool
actually batches into, and the sign disclosed in both directions (the model
also \emph{over}prices some shapes, minimum ratio 0.914). The $+21.8\%$ drift
appears only past $N{=}64$ --- beyond the FIFO cap, a regime the pool never
enters. Second, a cost-model error only matters if it moves a verdict. Third,
I demonstrated that it does not, rather than arguing it: the closed loop was
re-run with the measured vLLM cost substituted per dispatched batch --- p99
moved \SI{589}{\milli\second} to \SI{606}{\milli\second}, provisioned
worker-seconds identical to the digit, every verdict unchanged; and the
static \emph{denominator} was re-checked under measured cost too
(\SI{692}{\milli\second}, still feasible), so neither side of the ratio
carries a silicon penalty. The general principle I want on the record: the
right response to a model error is to \emph{price the conclusion under it},
not to defend the model.
\Point{Control/treatment pair: the control replay (ratio off) reproduces the
committed \SI{588.37}{\milli\second} baseline at \SI{589.30}{\milli\second};
the treatment (measured vLLM ratio on) gives \SI{605.61}{\milli\second} at
2{,}410{,}754.7 worker-seconds --- the same schedule, so the cost is the
schedule's, not the cost model's.}

\Q{The encoder measurements were once 1.4--1.6x off the model. What happened?}
\A That gap was real and I chased it down; it turned out not to be a model error
but an \emph{execution-engine} difference. The original measurement used
HuggingFace eager execution. When I measured the vision tower as the serving
engine actually runs it --- vLLM's kernels, which is what a deployment would use
--- the composed model holds within $12.8\%$ across the whole serving batch
regime. I did not quietly drop the old number: the $1.4$--$1.6\times$ is
disclosed and attributed, and the appendix carries the four-engine sweep.
\Trap{Do not present this as if I always knew. The honest story --- I found a
$1.5\times$ discrepancy, treated it as a real threat to the work, and ran it to
ground --- is stronger than pretending the model was never in doubt.}

\Q{Your reassurance that the AIC sizing ``carries no silicon penalty'' rests on 606 vs \SI{589}{\milli\second} from repricing the winning schedule with the measured vLLM cost. But you built that vLLM curve to be within 12.8\% of AIC, and you replayed a \emph{fixed} schedule --- so the reprice cannot move much and the worker-seconds are unchanged by construction. What did it falsify?}
\A Fair --- I will call it what it is: a feasibility check on a fixed schedule
under a near-identical curve, not an independent economic test, and
``worker-seconds unchanged'' is the same schedule replayed. The robustness
claim that carries weight is different: the same schedule under the
\emph{pessimistic} HuggingFace-eager cost (an upper bound on the engine gap)
stays feasible on every burst window (\SI{784}{} vs \SI{588}{\milli\second})
because the $\rho{=}15$/p90 headroom provisions 24--32 workers where even
that slower stack needs 12--16 --- so the engine-stack uncertainty is
\emph{absorbed by the forecast headroom}, with no added workers, and the
static $C{=}8$ denominator stays feasible under the reprice too (692 vs
\SI{647}{\milli\second}).
\Concede{Restate ``no silicon penalty'' as ``the 1.4--1.6$\times$ engine
uncertainty is absorbed by the forecast headroom.'' The 606-vs-589 number is
a feasibility check, not the evidence --- route the weight to the
headroom-absorption argument and the 12.8\% forward match.}

\Q{Does the encoder result survive on real silicon, or only in the cost model?}
\A It survives, and I tested it specifically \emph{because} the model underprices
the knee windows. Replaying the winning schedule with measured silicon cost
substituted: AIC p99 \SI{588}{\milli\second} becomes silicon p99
\SI{784}{\milli\second} --- both under the \SI{1}{\second} bound. The
p90/$\rho{=}15$ headroom provisions 24--32 workers at burst peaks against a
silicon-feasible $C^*$ of 12--16, so both hit their zero-queueing floors and the
underprice is absorbed with \emph{no added workers}. The schedule is unchanged
and feasible, so the cost ratio is preserved.
\Concede{The measured $f(N)$ used there is eager whole-tower --- a pessimistic
bound. A live encoder worker on the real request path remains future work, and I
should say that before being asked.}

\Q{State what each of your seven encoder-delay p99 values is --- 588, 589, 606, 644, 647, 692, 784\,ms --- and reconcile them. And explain why your ``${\approx}4\%$ replay noise, no quoted contrast within that band'' claim is not refuted by your own 606-vs-589 reprice, which is 2.9\%.}
\A They reconcile. \SI{588}{}: the zero-queueing service floor ($f(64)$ plus
the p99 handoff) and the burst-window forecast p99. \SI{589}{}: the forecast
week-scale aggregate under AIC. \SI{606}{}: that same schedule repriced at
measured vLLM cost. \SI{644}{}/\SI{647}{}: static $C{=}8$ aggregate under AIC
(the table rounds to the shared floor; the reprice sentence carries the raw
647). \SI{692}{}: static $C{=}8$ repriced at vLLM. \SI{784}{}: the forecast
schedule under pessimistic HuggingFace-eager cost. The ratio is stable
throughout: $2.41\text{M}/3.87\text{M}=0.62$.
\Concede{Two manuscript fixes. The ``no quoted contrast within 4\%''
sentence is self-refuting as written --- the 606-vs-589 reprice is
deliberately \emph{within} noise (that \emph{is} the non-divergence I want);
qualify it to ``no quoted \emph{policy} contrast.'' And print the raw
aggregates (589, 647) in the table rows rather than rounding to 588/644. The
606-vs-589 agreement shows non-divergence at replay noise, \emph{not}
engine-representativeness --- that weight belongs on the 12.8\% forward
match.}

\Q{What kind of model is the encoder cost model --- a roofline, a fitted analytic model, or a lookup table? Your extrapolation claims depend on which.}
\A On the regime the pool actually enters it is a \emph{lookup table of
measurements}, not a fit. AIC composes one encoder forward as an arithmetic
sum over the 27-block operator graph, one query per operator against the
collected tables (TP-communication terms analytic, zero at TP1). Each query
resolves in three strata: an exact-hit \emph{measurement} (mean of six
CUDA-graph replays), a piecewise-multilinear \emph{interpolation} between
bracketing collected points, or --- only past the grid edge, or only when a
system has no data at all --- an extrapolation or the pure-analytic
speed-of-light floor. After the shape-infill every 448-pixel attention and
GEMM query for $N\le64$ is an exact hit, so on its operational domain the
returned latency \emph{is} the silicon number. Roofline reasoning enters
only in the no-data case, which the 5090 collection eliminates for the
encoder.

\Q{The frames-class prefills that decide the burst TTFT tail are priced by extrapolation beyond your measured H100 grid. Is the headline burst result built on unmeasured data?}
\A About 3.1\% of trace requests (up to ${\approx}149$k tokens) exceed the
16{,}384-token context-attention grid and are priced by stratum-three
extrapolation on the H100 side --- and it is the frames-class prefills, so
I disclose it twice. Two mitigations. The extrapolation is not blind:
attention is analytically quadratic, and out to 147k-token contexts the
tables track that quadratic within a few percent per octave. More
decisively, the burst language verdict is a \emph{floor} argument --- a
${\approx}36$k-token prefill already busts the TP1 budget and no worker
count repairs it --- so whether 149k tokens are priced exactly or somewhat
off does not move the threshold crossing.
\Concede{This compounds the unvalidated-H100 point: the extrapolation is on
the language side. Exact 149k-token latencies validated on physical H100 are
honest future work.}

\Q{Your three silicon anchors validate the cost model. How do you know the \emph{simulator} is not introducing timing errors the anchors cannot see?}
\A I caught exactly such a bug and fixed it, which is the honest evidence
the harness is scrutinized rather than trusted. An early wheel-based
runtime priced each request through a per-request Python callback whose
\emph{own} wall-time leaked into the simulated clock --- it left p99
identical but inflated the extreme maximum (max TTFT 7.6 versus
\SI{4.5}{\second} on identical perf data). I diagnosed it as an
instrumentation artifact (a timing hook slower than the events it prices),
moved every calibrated run to the source-built runtime that evaluates AIC
predictions in-core in Rust in microseconds, and documented it. Then
determinism does the rest: same trace, same code path, same p99 to the
decimal --- the FIFO replay reproduces the simulated 54{,}169.5\,ms as
54{,}169.6.
\Point{Volunteer this: I found a simulator-side timing leak myself,
attributed it, and eliminated it. It shows the harness is audited.}

\Q{A 600\,W consumer card throttles under sustained load, and locking the SM clock does not guarantee it holds. Did you control thermal and power throttling?}
\A I locked SM to \SI{2407}{\mega\hertz} and memory to \SI{14001}{\mega\hertz}
(persistence mode) after finding an unlocked collection ran ${\approx}17\%$
fast --- isolated microbenchmarks ride higher clocks than sustained serving
--- and I sample the SM clock after \emph{every} replication, which is the
mechanism that would catch a throttle-induced sag below the lock. What I do
not separately report is board-power headroom or thermal steady-state. The
defense that matters: the validation \emph{ratios} cancel clock state,
scaling the collected table and the end-to-end measurement equally, so
residual throttling that hit both identically does not corrupt the
$3\%$/$12.8\%$ agreement.
\Concede{Absolute latencies could carry a throttle bias I did not isolate;
the relative validation is robust to it. Say it precisely.}

\Q{You justify the whole collection campaign by a ``$5.9\times$ optimistic'' speed-of-light fallback. Is that not a cross-device comparison dressed as a fallback error?}
\A The $5.9\times$ \emph{is} cross-device and I flag it in-line (SOL on the
5090's data-less die-sibling versus measured on the 5090 itself). But the
number is illustrative --- the load-bearing motivation is qualitative: on a
system without collected encoder data there is \emph{no soft fallback at
all}. SILICON mode raises \texttt{PerfDataNotAvailableError}; HYBRID and
EMPIRICAL raise \texttt{EmpiricalNotImplementedError}; only SOL answers,
from FLOP/bandwidth ceilings. Worse, VLMs show green in AIC's support
matrix on such systems only because the harness supplies no image
dimensions --- it silently skips the encoder phase and certifies text-only
costing. That masking is the real gap the collection closes.

\Q{Why an RTX~5090 for the encoder pool? That is not a datacenter part.}
\A Deliberately. The vision tower is 576M parameters and fits comfortably in
consumer memory, which makes consumer-class parts a plausible and cost-effective
hardware class for an \emph{independently scaled} encoder pool --- that is a
consequence of disaggregation, not an accident of what I had access to. It also
exposed a real gap: no encoder-timing data existed for \emph{any} consumer or
workstation GPU, and the planner has no usable fallback (its speed-of-light bound
is ${\approx}5.9\times$ optimistic). Collecting and upstreaming that data is
contribution~3.
\Trap{If pressed on ``you used it because it is what you had'' --- partially
true, and I should not die on this hill. The defensible core is that the gap in
the planner is real and the data is now contributed, regardless of motive.}

\Q{Why does the encoder's cost depend only on the \emph{total} images in a forward, whatever way they group into requests?}
\A Because the vision tower has no cross-image attention: each image is an
independent grid of patches, attention is evaluated at
batch$\times$images and every other operator at the total token count, so
the arithmetic is identical however $N$ is factored. I did not assume it ---
I probed AIC's graph directly, and $(1,64)$, $(2,32)$, $(64,1)$ return the
same latency \emph{to the microsecond}. That exact law is what licenses a
cross-request batching policy and lets each dispatched batch be costed at
its image total. Batching pays until it stops: $f(1){=}\SI{8.16}{\milli\second}$
per image, $f(64){=}\SI{316.4}{\milli\second}$ (\SI{4.94}{\milli\second}
amortized, a $1.65\times$ gain over a ${\sim}\SI{3.2}{\milli\second}$ fixed
per-forward overhead), and the realized benefit saturates by
$N{\approx}6$--$8$ on both execution stacks --- a compute ceiling that
bounds the useful dispatch batch, not a cost-model artifact.

\Q{The inter-pool handoff is an analytic bandwidth bound. Is that not the weakest-modeled part of the pipeline?}
\A It is the most simplified part, and its two biases run in opposite
directions, which I state. The bound is $\text{RTT}+\text{bytes}/\text{bandwidth}$
on a parameterized link (default \SI{25}{Gbps}) for the $196\times4096$
bf16 embedding stream (${\approx}\SI{1.6}{\mega\byte}$ per image) ---
bandwidth-only, ignoring contention, serialization, NIC/PCIe placement and
cross-node topology, so it is a \emph{lower} bound on transfer latency.
Against that, transfer is charged to request latency but \emph{not} to the
worker (asynchronous DMA), which is optimistic for encoder throughput. It
is deliberately replaceable: a measured NIXL transfer curve drops in
without touching anything downstream. And I already price the two known
deltas --- the DeepStack $4\times$ payload (the fabric-requirement
limitation) and the burst-window p99 of \SI{272}{\milli\second} that sets
the \SI{588}{\milli\second} floor.
\Concede{No measured interconnect curve --- a bandwidth-only bound. Say it,
and say a NIXL measurement is the clean next step.}

\Q{Chunk-parallel dispatch conveniently makes encoder tail latency capacity-controllable. Is that not a self-serving simulator assumption?}
\A It is a simulator policy assumption and I flag it as one --- but it is
physically grounded, and I show the result does not hinge on it. Images are
independent in a non-causal encoder, so splitting a $97^2$-image request
into ${\le}64$-image forwards is what an encoder-disaggregated runtime does;
it is a property of the policy, not of the trace. I keep the sequential
alternative behind a flag --- it produces an irremovable ${\approx}2.9$\,s
p99 floor, flat from 4 to 32 workers --- and I show the cap does not drive
the burst verdict (p99 5.2/5.0/\SI{6.7}{\second} at caps 32/64/128). The
vision-pool controller is my own event-driven FIFO inside the converter
(Dynamo's mocker models text tokens only); the language pool is Dynamo's
real mocker/replay with a KV-aware router.

\Q{Do you model KV-cache pressure, eviction, and admission control in the language pool?}
\A No. The language pool inherits exactly what Dynamo's mocker and
\texttt{dynamo.replay} model, and I claim no KV eviction, preemption, or
admission control on top. What is modeled is prefix-cache reuse via the
KV-aware router (disjoint prefix-hash sets, so no spurious cross-image
reuse), with prefill length taken verbatim from the trace's
\texttt{ContextTokens}. The replay is open-loop: every request is admitted
and priced by its prefill/decode cost; ingress (image fetch, preprocessing,
tokenization, front-end queueing) is outside the boundary. This is
conservative in one direction (no admission control lets the queue build,
which is what exposes the episodic blocking) and optimistic in another (a
real system hitting KV limits would preempt or reject, possibly capping the
very tails I measure). The burst failures are prefill-compute floors and
encoder FIFO blocking --- arithmetic and queueing phenomena robust to
KV-modeling detail.

\Q{Is the reported p99 even statistically meaningful over these windows?}
\A Yes, by volume and by construction. The burst windows contain tens of
thousands to hundreds of thousands of requests (burst1 alone spans 64{,}804
against the encoder pool; the week-scale spans 749{,}601), so the p99
estimator sits on hundreds to thousands of tail observations, not a handful.
Determinism does the rest: the simulation is open-loop and seeded --- the same
trace, same arrivals, same code path reproduce the same p99 to the decimal
(the independent multi-worker silicon harness re-derives the committed grid
optimum for burst3, $C^*{=}7$ at \SI{904.7}{\milli\second}, exactly; the
end-to-end FIFO replay reproduces the simulated p99 as 54{,}169.6 versus
54{,}169.5\,ms), so cross-policy deltas are not fighting sampling noise. Where a stochastic
component does exist --- the live planner-in-the-loop runs --- I measured
run-to-run spread (${\approx}4\%$ on p99), report single-run differences as
point differences rather than noise bands, and built a mandatory
${\approx}14\%$ capacity margin into the tail-aware conversion precisely
because of it.
\Trap{Do not say ``the simulator is deterministic so there is no variance
question.'' The honest form: determinism removes \emph{sampling} noise from
policy comparisons; the residual question is \emph{model} error, and that is
answered by the silicon anchors, not by repetition.}

\Q{Ten rounds of review found real errors. Why should I trust the numbers now?}
\A You should trust them partly \emph{because} of that --- but let me describe
the process rather than ask for credit. Every number in the document is generated
by a versioned script over archived inputs; every result carries a ledger entry
naming the artifact, the exact command, and which paragraph it feeds; a number in
a chapter with no ledger entry is treated as a defect. The audits found real
errors --- an artifact-inflated denominator, an actuation-timing bug, an
over-claimed mechanism --- and each is recorded with what it superseded. What I
would ask you to notice is the \emph{direction} of the corrections: they all moved
the headline against me, from $0.27\times$ to $0.62\times$, and the descriptive
core --- the decoupling, the forecastability, the silicon cost model, every SLA
verdict --- did not move at all.
\Trap{Tone decides this one. It reads as either defensive or confident depending
entirely on delivery. Name the errors \emph{first}; earn the ``but the core
survived'' clause rather than rushing to it.}

\Q{Walk me through re-running the 0.62$\times$ from a clean checkout. The AIC data is an unmerged draft PR; the aiconfigurator and dynamo venvs are external with no dependency lock; a live-planner fix is a local-only commit; some scripts hard-code paths; and you ran the replay on two runtimes that disagree in the tails.}
\A Two claims, and I keep them apart. \textbf{(1)~The headline is
artifact-exact.} From a detached clean worktree at HEAD, with no live-tree
inputs, the runner recomputes all 313 comparator cells from the tracked raw
artifacts --- zero missing statistics, zero mismatches --- and the predicate
lands at $0.623$ recomputed against $0.623$ expected, with the composed
nine-cell artifact regenerated byte-identically to the committed object
(\texttt{verification/repro\_clean\_checkout\_report.md}; path and
byte-identity requirement defined in
\texttt{supplementary/SUPPLEMENTARY.md}). The reproduction spine underneath
it is real: the headline operands discussed here map to versioned scripts
and archived artifacts in the manifest, the ledger records
artifact/command/date/feeds/supersedes, and the headline runtime is
unambiguous --- the 1.3.0 source build (in-core Rust predictions) produced
every calibrated run from the load sweeps on, and the 1.2.1 wheel's
Python-callback leak inflated only the extreme max, with identical p99. The
load-bearing issue-time pipeline resolves paths from \texttt{TESE\_ROOT} and
relocates. \textbf{(2)~A rerun of the original platform is not yet fully
publicly reproducible.} Regenerating the external
aiconfigurator/Dynamo environments and the live anchor from public artifacts
is a different claim: those two venvs are not fully locked and the live-planner
fix is a local-only commit.
\Concede{Concede the second gap plainly --- the external environments and the
local fix are not fully pinned, so the live anchor cannot be rebuilt end to
end from public artifacts today. Do not let that concession travel to the
first claim: the recomputation of the headline from tracked raw artifacts is
exact, not within-noise.}

\Q{You contributed a Dynamo fix (\texttt{fpm\_observation\_mode}) that makes the live planner ready at its first tick. Did you patch the platform to make your policy pass?}
\A No --- the fix is policy-neutral, and its absence would have hurt the
comparison's \emph{fidelity}, not favored either side. Live, Dynamo's
planner samples only the latest forward pass per worker per tick, and the
mixed prefill+decode passes that inform aggregated capacity are rare at calm
load (zero arrived in five minutes at 2\,req/s), leaving the capacity model
unready until the workload was already busy --- too late for \emph{any}
scale-ahead policy, forecast or reactive. The opt-in fix drains every
published pass into the same regression feed the offline replay already
used, so the live planner matches its offline twin's first decision. It
restores parity between the two harnesses; it does not tilt the
forecast-versus-reactive verdict, which reproduces in both directions
regardless.
\Point{The demand path is byte-identical live and offline (first tick
\texttt{29.86 rps}, 6 replicas in both). The fix is on the capacity-belief
warmup, symmetric across policies, not on the forecast.}

\Q{Why TP2 for an 8B model? TP1 replicas are the textbook answer.}
\A Because the burst bottleneck is a per-request \emph{prefill-latency floor},
not throughput --- and only intra-request tensor parallelism shortens a single
prefill that independent replicas can only queue behind. The frames class
produces ${\approx}36$k-token prefills; at iso-cost --- four H100s either way
--- TP1$\times$4 fails every burst window (TTFT p99 1{,}314--1{,}331\,ms)
while TP2$\times$2 passes every one (763--846\,ms), a $1.55$--$1.74\times$
tail reduction from the parallelism \emph{layout} alone. And the floor is
structural, not budgetary: at TP1 \emph{no} replica count is feasible ---
the long prefill pins TTFT p99 at ${\approx}1{,}285$\,ms even at $N{=}16$,
sixteen GPUs failing to buy what four buy at TP2. Replica count buys
throughput; layout buys the floor; this workload's binding constraint is the
floor.
\Point{Table \texttt{base-tp}: the clean iso-cost A/B across all five burst
windows. This is also a transferable design rule: long-prefill multimodal
traffic changes the optimal parallelism layout, not just the fleet size.}

\Q{Why build your own two-pool simulator instead of using an established one --- Vidur, LLMServingSim?}
\A Because none of them models the stage this thesis is about. Vidur and
LLMServingSim are text-only LLM-inference simulators (prefill and decode)
with no vision-encoder pool, and Dynamo's own mocker --- which I \emph{do}
use for the language pool --- models text-token processing only, which is
exactly why the encoder pool cannot live inside it. I did not build a new
LLM simulator: I composed an AIC-costed encoder FIFO (with its
probe-verified batching law and analytic handoff) with Dynamo's production
\texttt{dynamo.replay} through the trace-timestamp seam, so the language
pool is a real production control path and the vision pool is
silicon-costed. The contribution is the composition and the per-pool
attribution, not a from-scratch engine --- I cite Vidur as the
established-instrument precedent rather than reinvent it.

% ============================================================================
\section{Theory, definitions, and terms}
\label{sec:theory}

Not questions, but the vocabulary the questions are built from. A committee
member may ask ``what exactly do you mean by X'' at any point, and fumbling a
definition costs more credibility than conceding a limitation. Each entry
below is written the way the thesis itself uses the term, with the place in
the manuscript and the measured instance attached, so the answer is a
definition \emph{plus} a number rather than a textbook recitation. The third
part of this material, the theory quick-facts card, sits in the cram layer
(Section~\ref{sec:cram}) beside the rapid-fire numbers so that both can be
drilled together.

\subsection*{Definitions as this thesis uses them}
\addcontentsline{toc}{subsection}{Definitions as this thesis uses them}

\paragraph{SLO versus SLA.} An SLO is the general concept: a service-level
\emph{objective} on serving quality, conventionally stated on upper-tail
percentiles. This document reserves ``the \emph{evaluation} SLA'' for the
specific tail bounds every policy is graded against. That convention is
declared in the introduction (\texttt{tex/intro.tex}, opening section) and
followed thesis-wide. The evaluation SLA is three bounds: encoder-delay p99
$<\SI{1}{\second}$, TTFT p99 $<\SI{1}{\second}$, ITL p99
$<\SI{100}{\milli\second}$. Their provenance is $2\times$ Dynamo's own
planner defaults (\SI{500}{}/\SI{50}{\milli\second}), which the planner
applies as window-\emph{means} and which this thesis applies as \emph{tails}.
If asked whether SLA is the right word: it is a naming convention borrowed
from serving-systems usage, not a contractual claim.

\paragraph{TTFT.} Time to first token, dominated by queueing and prefill. In
the two-pool system an image-bearing request's first token waits on encoder
delay \emph{and} then on prefill, which is why the two pools are graded
separately and never summed into one bound.

\paragraph{TPOT versus ITL.} TPOT is time per output token, the decode-phase
statistic. This thesis reports it throughout as the equivalent per-gap
inter-token latency, ITL, defined in \texttt{tex/metodologia.tex} as ``the
per-token decode gap that TPOT summarizes''; in the live client the two are
the same quantity ($\text{ITL}=\text{TPOT}$). So ``ITL p99
$<\SI{100}{\milli\second}$'' is the same contract as a per-token decode
bound, read per gap rather than per request.

\paragraph{Percentile, quantile, and which estimator.} A quantile is the
value below which a given fraction of the sample lies; a percentile is the
same thing named on a 0--100 scale. Every reported p99 in this thesis is an
\emph{order-statistic selection with no interpolation}: sort the sample and
take index $\min(n-1, \lfloor n p/100 \rfloor)$. That is the converter's own
rule (the \texttt{percentiles} helper in
\texttt{workload/trace\_to\_two\_pool.py}) and it is reused verbatim by the
independent recompute (\texttt{pctl} in the comparator frontier's
\texttt{build\_frontier\_summary.py}) and by the metric stack. Two consequences worth knowing cold.
Every quoted p99 is an \emph{observed} request delay, never a weighted blend
of two neighbours, so it is reproducible to the decimal under determinism.
And it is a rank-selection estimator in the nearest-rank family, not the
linear-interpolation default of \texttt{numpy.percentile}, which is why arms
sitting on the same service plateau report the \emph{identical} value
(\SI{589.301}{\milli\second}) rather than nearly-identical ones.

\paragraph{The p90 band edge, and what ``calibrated'' means for it.}
Coverage is the fraction of realized values at or below a predicted
quantile. A p90 edge is calibrated if that fraction is $0.90$. Measured here:
$0.88$--$0.90$ against nominal $0.90$ for every model and signal at both
provisioning cadences, the widest deviation over the twenty combinations
being $0.880$. Resolved by horizon step on the image signal it is
$0.876$--$0.903$ at \SI{15}{\second} bins and $0.857$--$0.911$ at
\SI{1}{\minute} bins with no trend in $h$; split by burst condition it is
$0.890$--$0.901$ in burst against $0.890$--$0.896$ outside. The lower edge is
\emph{not} usable: the p10--p90 interval covers 69\% against a nominal 80\%,
broken by zero ties, which is exactly why provisioning is one-sided on the
upper edge. Say ``a calibrated upper quantile,'' never ``a calibrated
predictive distribution.''

\paragraph{WAPE.} Weighted absolute percentage error,
$\text{WAPE} = \sum_t |\hat y_t - y_t| \big/ \sum_t y_t$, as printed in
\texttt{tex/predictors.tex} and implemented in
\texttt{predictors/additive\_structural\_ablation.py}. It is scale-free and
defined on zero-signal bins, which MASE's in-sample-naive denominator is not
on a heavily zero-inflated, non-stationary series. Its no-skill floor is
100\%: predicting identically zero scores exactly 1.0. That floor is the
reason the \SI{1}{\second}-bin image number (98.8\%) must never be quoted
without it.

\paragraph{Quantile (pinball) loss and CRPS.} The proper scoring rules for a
probabilistic forecast: pinball loss scores one quantile, CRPS integrates
over all of them. \emph{Neither is computed in this thesis} --- the correct
answer is that probabilistic quality is assessed instead by empirical
coverage of the single edge the controller consumes, plus a direct exceedance
check. Concede the omission; do not invent a number for it.

\paragraph{Rolling-origin evaluation.} An origin $o$ predicts bins
$[o, o{+}H)$ from context $[o{-}C, o)$, with $H{=}12$ bins at the series'
native cadence, and origins advance in strides of $H$ so the forecasts tile
the timeline into a continuous, non-overlapping path. That is exactly how a
scaler consumes forecasts: stand at a point in time, ask for the next $h$
steps, act, let the clock advance. Everything is scored on the \emph{common
span}, the origins predictable by every model, from 167 per model at
\SI{1}{\minute} up to \num{49727} at \SI{1}{\second}. The common-span
threshold is an \emph{index}, not a wall-clock time, so it covers 98.7\% of
the week at \SI{1}{\second} bins but only the final \SI{33.4}{\hour} at
\SI{1}{\minute} bins. Within-column comparisons are unaffected; cross-column
WAPE levels are not comparable.

\paragraph{Paired-origin bootstrap.} A resampling test on the \emph{matched
per-origin} score differences between two forecasters: resample origins with
replacement, recompute the gap, read a 95\% interval. The thesis's is stated
as ``300 independently resampled time-ordered origins, not a day- or
block-level population test'' --- the qualifier matters, because it bounds
the test to origin-level resampling and disclaims a population claim. It is
used twice: the sub-point ARIMA-versus-Chronos gap at \SI{15}{\second} is
\emph{not} distinguishable from zero, while Toto with covariates at 60.9\%
beats ARIMA's 66.8\% with an interval that excludes zero.

\paragraph{Zero-shot forecasting.} A pretrained model applied to a series it
was never fitted on: no per-signal training, no fine-tuning, no trace
history, no hyperparameter search. It is the \emph{conservative} reading of a
foundation model's capability, being what an operator gets on day one. All
three foundation models here run zero-shot, and the headline rests on the
smallest of them (Chronos-2, 120M).

\paragraph{Issued lead versus forecast target horizon.} These are not the
same axis, and confusing them cost this thesis a factor of two. The
\emph{target horizon} $h$ is how far ahead of its origin a forecast bin sits:
$h=1\ldots12$, i.e.\ \SIrange{15}{180}{\second} ahead at \SI{15}{\second}
bins. The \emph{issued lead} is the wall-clock time between the forecast
being available and the bin it describes: issue time $=$ target
$-\,(h{-}1)\cdot\SI{15}{\second}$. So $h{=}1$ is a nowcast carrying
\emph{zero} usable lead, and an $H{=}12$ archive at \SI{15}{\second} bins
supplies at most $(12{-}1)\times 15 = \SI{165}{\second}$ of issued lead. The
closed loop's declared grid is therefore $\{15,30,45,165\}$\,\si{\second}.
The old \SI{300}{\second} cell was target-indexed: an evaluation bound, not
an issued lead.

\paragraph{Issue-time versus target-time indexing.} Issue-time actuation
means every scaling decision uses only forecasts that had actually been
issued by the moment it acts. Target-time indexing dates a forecast to the
bin it predicts and lets the schedule read it as soon as that bin arrives,
which is an optimistic upper bound no issue-time schedule realizes. Rebuilding
on honest issue-time actuation is one of the two corrections that moved the
headline from $0.27$--$0.31\times$ to $0.62\times$.

\paragraph{Readiness models: warm-pool versus strict re-warm.} These describe
what happens when a scaled-down worker is wanted again. Under \emph{warm-pool}
readiness a scaled-down replica returns without re-paying cold start. Under
\emph{strict re-warm} it pays the full cold-start lead again. The shipped
mechanics are asymmetric --- the forecast and oracle arms actuate through a
warm pool while the reactive family re-warms strictly --- and both directions
were repaired and measured rather than argued. One caveat is disclosed and I
must repeat it: the warm-pool direction is algorithmically symmetric, but the
forecast's strict re-warm variant retains matured orders as reusable for a
\SI{165}{\second} hold, close to though not identical with the reactive arm's
trailing-window re-warm semantics, so that endpoint is a sensitivity, not a
fully equalized model.

\paragraph{The threshold-integrated identity.} Total delay mass is exactly
the area under the complete empirical exceedance-count curve:
$\sum_i d_i = \int_0^\infty N(d_i > t)\,\mathrm{d}t$. That is why it is the
convention-minimized leading statistic --- it integrates \emph{every}
possible threshold instead of privileging one, so it introduces no SLA cut
and no percentile selection of its own. It is not convention-free: it stays
conditional on the declared evaluation population and span, on the
encoder-delay definition, and on a \emph{linear} loss functional. A quadratic
or capped loss could rank arms differently, and no such ranking is computed.
It is a post-hoc synthesis that may lead the prose but does not replace the
predeclared every-window compliance endpoint.

\paragraph{Excess mass.} $\sum(\text{delay} - \SI{1}{\second})^{+}$ over the
evaluation population, in request-seconds: not how many requests break the
bound but how far past it they sit. Over the aggregate span the forecast arm
carries \num{3688.0} against reactive's \num{6709.7} (45.0\% less), with
static $C{=}8$ at \num{289.9}. It is the statistic that separates ``a few
requests slightly late'' from ``a queue that diverged.''

\paragraph{Masking and anti-masking.} Masking is an aggregate statistic
hiding a localized failure. The live instance inside the evaluation itself:
the reactive rows \emph{appear} to pass because their aggregate p99 clears
the bound, while they fail the fastest onset outright, because that window's
${\approx}3{,}100$ requests are under half a percent of the span. That is the
same dashboard-masking the trace chapter finds operationally (a pool
violating its tail bound at 29\% mean utilization while every
minute-averaged series reads healthy), recurring one level up inside the
metric. Anti-masking is the response: grade aggregate \emph{and} per-window
jointly, never either alone.

\paragraph{The every-window compliance gate, and its dual role.} The gate is
the predeclared, frozen endpoint: every declared evaluation window must hold
its p99 bound. It does two jobs at once, and the thesis says so explicitly.
First, it is the \emph{compliance} contract, the thing a policy must pass to
be called passing. Second, it \emph{doubles as the masking check} on the
aggregate statistics: it shows where the tail is lost and certifies that no
aggregate hides a localized failure. The p50 row is the live demonstration of
why that certificate is kept --- an excellent cost policy that fails the SLA
outright.

\subsection*{Theory the results lean on}
\addcontentsline{toc}{subsection}{Theory the results lean on}

\paragraph{Long-range dependence.} A process whose autocorrelation decays
slowly, as a power law rather than exponentially, so correlations survive at
lags far beyond any natural timescale and the variance does not smooth out
under aggregation. This is the property that makes threshold rules on
averaged signals unsafe: widen the averaging window and the variance grows
faster than the mean. Measured instance: the image signal's autocorrelation
half-life is \SI{69}{\minute} against \SI{9}{\minute} for the request signal
at the same one-minute resolution, on a $1/f^{0.86}$ spectral background. The
operational consequence, and the only one the thesis needs, is that no tested
averaging window makes the signal safe for a reactive threshold.

\paragraph{Hurst exponent.} The self-similarity parameter: $H = 0.5$ is a
memoryless process, $H > 0.5$ long memory, and $H \to 1$ is strong
persistence. Measured instance: $H \approx 0.93$ for images against
${\approx}0.83$ for requests. The honest scope, printed in the manuscript, is
that $H$ is derived algebraically from the spectral slope,
$H = (1{+}\beta)/2$, over the sub-diurnal band (10 minutes to 12 hours). It
is one estimator viewed two ways, not two corroborating measurements: no
R/S, DFA, or wavelet estimator was run, and the fit is endpoint-sensitive
(refitting with the band's upper period at 6 hours gives $H \approx 0.90$, at
3 hours $H \approx 0.82$). Defend the images-versus-requests \emph{contrast},
not the constant.

\paragraph{Autocorrelation half-life.} The lag at which the autocorrelation
function falls to half its value at lag zero: a model-free persistence
statistic needing no self-similar model at all. Measured instance across the
four demand signals: images \SI{69}{\minute}, context tokens
\SI{15}{\minute}, requests \SI{9}{\minute}, generated tokens
\SI{5}{\minute}. It is \emph{not} deseasonalized, which is why it is used
comparatively only. Half-life and Hurst orderings disagree across the four
signals, the decode series pairing the shortest half-life with the
second-highest $H$, so the thesis makes no joint claim on them.

\paragraph{Self-similarity.} Burstiness that looks statistically the same at
every timescale, so aggregation does not smooth it. The point of invoking it
is positional rather than technical: this is the \emph{normal} finding for
aggregated human demand, established for Ethernet traffic by Leland et
al.~(1994) and for Web traffic by Crovella \& Bestavros~(1997), so a serving
trace showing $H \approx 0.93$ is continuous with thirty years of
teletraffic measurement rather than an outlier that needs luck to reproduce.

\paragraph{The Taqqu--Willinger heavy-tailed ON/OFF duality.} Superposing
many ON/OFF sources whose ON periods are heavy-tailed produces long-range
dependence in the aggregate: heavy tails and long memory are two readings of
one mechanism, not two properties that coincidentally co-occur. That is the
correct framing for this trace, and I should volunteer it rather than sell a
surprising coupling. The frames class (2.7\% of requests, 81.2\% of image
volume) drives \emph{both} the second moment (CV$^2$ 3.83) and the
low-frequency spectral power ($H \approx 0.93$). What the operational claim
needs is not that the two are independent, but that the envelope is
\emph{forecastable} at provisioning leads, and that is measured directly
rather than inferred from the duality.

\paragraph{Fano factor, index of dispersion, and overdispersion.} The Fano
factor of a counting process is $\operatorname{Var}(N_T)/\mathbb{E}(N_T)$ for
counts in windows of width $T$. A Poisson process has Fano $=1$ at
\emph{every} $T$; a value above 1 is overdispersion, and Fano growing with
$T$ is the counting-process signature of long-range dependence. Measured
instance: 1.8 at \SI{1}{\second} growing to 938 at one hour, with
inter-arrival coefficient of variation 1.55 and per-second counts
overdispersed against Poisson$(1.65)$. Every aggregation scale an autoscaler
would plausibly use is therefore overdispersed. The scope I must keep: a
seven-day record cannot establish asymptotically unbounded growth, and the
manuscript says so; the load-bearing claim is only \emph{within-window}
overdispersion, which is directly measured and is all a threshold rule needs
to be unsafe.

\paragraph{CV$^2$ as a burstiness measure.} The squared coefficient of
variation of the binned rate, $\operatorname{Var}/\text{mean}^2$: a
dimensionless, scale-free burstiness index. Measured instance at one-minute
bins: image rate 3.83 against request rate 0.32, an order of magnitude apart
on the same clock, with peak-to-mean $15.6\times$ against $4.7\times$. This
is the cleanest single number for ``the vision signal is a different kind of
signal,'' and unlike $H$ it needs no estimator argument.

\paragraph{Spearman rank correlation.} Pearson correlation computed on ranks,
so it detects monotone coupling whether or not the relationship is linear.
Measured instance: image count against generation length at request level,
$\rho = -0.06$, with no monotonic trend in per-bucket median generation. The
certified scope is ``weakly rank-correlated,'' and it is one of \emph{three}
tested relationships, alongside constant-fraction scaling and lagged linear
cross-correlation. The manuscript explicitly does not rule out nonlinear or
multivariate coupling on these tests; that direction is probed by the
covariate experiment, which finds nearly nothing to exploit at provisioning
cadence.

\paragraph{Little's law and queueing intuition.} Little's law states
$L = \lambda W$: mean number in system equals arrival rate times mean time in
system. It is true for any stable system and assumes nothing about the
service distribution. \emph{I do not invoke it, and I should not.} There is
no arrival model, no service-time distribution, no traffic intensity, and no
stability analysis in this work, and reaching for queueing machinery invites
a demand for all four. The thesis's own explanation of why p99 explodes at
29\% mean utilization is a deterministic one, and it is stronger: mean
utilization is the time-average of a \emph{bimodal} occupancy. The pool
drains fully between excursions and saturates inside them, so the average
measures work done while the p99 is set by the excursions. Inside the burst
window the backlog exceeds 1{,}000 queued images for 805 seconds, peaking at
4{,}701, with the worst second landing 2{,}726 images at once; encoder-delay
p99 reaches \SI{5.0}{\second}, $5\times$ the bound, while the one-minute
image-driven busy fraction reaches 0.90 and the minute-mean stays low. Call
it a deterministic event-driven FIFO replay exhibiting head-of-line blocking
that a mean monitor misses. Never say M/M/c, M/G/1, or PASTA.

\paragraph{Saturation knee.} Defined operationally in this thesis, not
asymptotically: the largest swept arrival speedup $s$ at which a pool still
meets \emph{every} SLA bound, reported together with that knee's offered
load. Measured instance on the steady head: the vision pool is uncensored
through $s{=}32$ ($\geq 256$\,img/s) while the language knee sits at $s{=}2$
under calibrated timing (5.5\,req/s) or $s{=}3$ under polynomial timing
(8.3\,req/s), a factor above ten apart, and at the calibrated language knee
the encoder idles at 4.5\% utilization. Both language knees are
\emph{ITL-bound}: the successor point still has $1.8$--$2.1\times$ TTFT slack
while the inter-token bound has already crossed, so that pool's capacity
under an interactive SLA is set by sustained per-user token throughput. On
the burst window neither pool has a knee: both are already over their bounds
at $s{=}1$.

\paragraph{Tensor parallelism, and what TP2$\times$2 means.} Tensor
parallelism splits each layer's weights across GPUs so that one forward pass
is executed jointly, which shortens a \emph{single} request's prefill.
Independent replicas cannot do that; they can only let a long prefill queue
behind itself. The thesis's notation, declared in the introduction, is
explicit: TP2$\times$2 means \emph{two tensor-parallel-2 replica groups}, so
four H100s in total, the same iso-cost budget as TP1$\times$4. Measured
instance: at that iso-cost budget TP1$\times$4 fails every burst window while
TP2$\times$2 passes all five at 763--\SI{846}{\milli\second}, a
$1.55$--$1.74\times$ tail reduction from the parallelism \emph{layout} alone,
and at TP1 no replica count is feasible at all --- TTFT p99 stays pinned near
$1{,}285$\,ms even at $N{=}16$. Replica count buys throughput; layout buys
the floor.

\paragraph{Prefill versus decode compute character.} An LLM answers in two
phases of opposite character: prefill is a single pass over the whole prompt,
compute-bound and parallel; decode emits one token at a time and is bound by
memory bandwidth. Service quality is specified against that anatomy, TTFT
against queueing plus prefill and ITL against decode. Two measured
consequences to have ready. The burst failure on the language side is a
\emph{prefill} floor, a per-request compute cost that no worker count
removes. The steady-state language limit is a \emph{decode} limit, which is
why the knee is ITL-bound rather than TTFT-bound.

\paragraph{KV cache.} The per-request cache of attention keys and values that
makes decode incremental instead of quadratic. It matters here twice. First
by absence: the vision encoder's attention is non-causal and carries no KV
cache, which is one of the three reasons the encoder stage resembles neither
prefill nor decode (the others being that its cost scales with image count
and resolution rather than conversation length, and that its work arrives in
bulk). Second by scope: in the language pool what is modeled is prefix-cache
reuse through the KV-aware router, with disjoint prefix-hash sets so there is
no spurious cross-image reuse. No KV eviction, preemption, or admission
control is modeled, and none is claimed.

\paragraph{ViT patching arithmetic for this tower.} The validation target is
the vision encoder of \texttt{Qwen/Qwen3-VL-8B-Instruct}: a 27-block ViT,
hidden size 1152, 16 heads of dimension 72, GELU MLP of width 4304, a
PatchMerger projector into the language model's 4096-dimensional embedding
space, and three DeepStack mergers. The arithmetic to have ready: at patch
size 16 a $448^2$ image is a $28\times28$ grid, hence \emph{784 patches} ---
the number the collection grid had to be made exact at --- and the
$2\times2$ PatchMerger reduces those to $784/4 = 196$ post-merge
tokens per image. That 196 is what the seam ships: a $196 \times 4096$ bf16
stream, ${\approx}\SI{1.6}{\mega\byte}$ per image. The disclosed mismatch is
that the language pool's prefill instead uses the trace's
\texttt{ContextTokens} verbatim, which embeds the \emph{traced} model's
${\approx}66$ tokens per image, roughly $3\times$ below a Qwen3-VL-consistent
196-per-image prefill.

\paragraph{Modality disaggregation.} Running the vision encoder as its own
pool, on its own hardware class, scaled independently of the language model.
The thesis adopts that system model from ModServe and extends it on the
\emph{control} side; the provisioning question only exists once the pools are
separate. Two consequences I should state without being asked. Mixing
hardware classes is the architecture, not an accident, so the encoder is
priced in RTX-5090-class worker-seconds and the language pool in H100
GPU-hours, and the two are never summed. And the pools bill from different
instants: language GPU-hours bill from instance \emph{order}, the encoder
pool from \emph{readiness} with its cold start reported separately as a
warm-up column.

\paragraph{Cold start anatomy.} A scale-up issued at $t$ delivers capacity at
$t + L$, and $L$ decomposes differently per pool. For the encoder the
measured \emph{model-side} minimum is \SIrange{3}{5}{\second}, and the tower
is small (576M parameters), so weight loading is not the binding term:
anything above ${\approx}\SI{5}{\second}$ is \emph{orchestration}-bound ---
pod scheduling, image pull, node acquisition --- which weight-multicast
advances do not touch. For the language pool the lead is the weight-loading
regime, tens of seconds and up, and the runs sweep $L \in
\{60,\ldots,300\}$\,\si{\second}. The encoder comparator grid sweeps $L \in
\{15,30,45,300\}$\,\si{\second}, above that measured model-side minimum,
while the forecast arm's declared grid is $\{15,30,45,165\}$\,\si{\second}
because that is the largest lead its archive can \emph{issue}.
% ==========================================================================
\section{Scope, limitations, and prepared concessions}
\label{sec:scope}

The boundaries I draw on purpose, and the ones I would rather volunteer than have excavated. Conceding a real limitation quickly is almost always cheaper than defending it badly, and every item here has a decided-in-advance concession attached.

\Q{The decode signal is an offered-work upper bound. Does that not undermine RQ1?}
\A It bounds one half of RQ1, and I flag it everywhere the decode signal appears.
The trace stamps each request's full generated-token volume at its \emph{arrival}
minute, so the decode series is offered work at arrival, not instantaneous decode
occupancy. That is faithful for encoder and prefill work --- where the thesis's
load-bearing claims live --- but it makes the decode correlation
($R^2 = \num{0.654}$) an upper-bound reading. Note the direction: this cuts
\emph{against} my own convenience. The decode signal is the one that looks most
rate-predictable, so if anything the true decoupling is at least as bad as I
report.
\Point{The limitation weakens the ``largely yes for decode'' half of RQ1 --- not
the vision half, which is the half the thesis rests on.}

\Q{The inter-pool seam serializes encoder and prefill. Is that not artificially pessimistic?}
\A It is conservative, deliberately. There is no speculative overlap: a request's
prefill does not begin until its encoder work completes. A real system could
overlap text-token prefill with image encoding and achieve a lower TTFT. That
makes my saturation knees \emph{pessimistic} about where the pools break --- which
strengthens rather than weakens an infeasibility claim. I would not want to be
caught claiming the assumption is neutral; it is not, it is favorable to my
critics.

\Q{The handoff assumes an embedding-only payload. What happens at full DeepStack?}
\A It breaks --- and it breaks for \emph{every} policy, including the static
baseline, which is why I present it as a fabric requirement rather than a policy
weakness. At Qwen3-VL's full DeepStack payload ($4\times$) on the modeled
\SI{25}{Gbps} link, handoff p99 goes from 272 to \SI{1087}{\milli\second} and
transfer alone violates the bound --- static $C{=}16$ included, at
\SI{1311}{\milli\second}. The \SI{588}{\milli\second} floor and the $0.62\times$
therefore assume an embedding-only handoff or a ${\geq}\SI{100}{Gbps}$ fabric. I
state this in scope, not in a footnote.

\Q{The seam ships 196 tokens per image, but the language prefill uses the trace's \texttt{ContextTokens} verbatim. Does that inconsistency invalidate the burst-window language verdict?}
\A It is the chapter's most surgical limitation, and I engage it directly
because the biases are aligned and disclosed. The two conventions describe
the same tensor two ways: a Qwen3-VL-consistent frames prefill would be
${\geq}529\times196\approx104$k tokens against the trace's recorded
${\approx}36$k, so frames-class language work is \emph{under}counted (making
the language pool look better in bursts), while the uniform $448^2$
assumption likely \emph{over}states frames-class encoder work (making the
vision pool look worse). Both biases lean toward the burst asymmetry I
report. And I state exactly which verdicts survive: the TP1 prefill floor is
robust (at $3\times$ the tokens it only worsens), the vision verdict is
convention-independent, but the TP2 burst-window \emph{pass} is
convention-dependent --- under the 196-per-image convention the frames
prefills would strain even TP2.
\Trap{Do not defend the TP2-passes-the-window claim as convention-robust;
the thesis itself marks it convention-dependent, and a both-conventions run
is future work. Defend the floor and the vision verdict; concede the TP2
margin.}

\Q{Everything is measured on the 8B. Does the result hold for the flagship 235B?}
\A The vision result transfers unchanged; the language envelope shrinks,
which only sharpens the thesis. The 235B shares the 8B's vision tower
\emph{and} its 4096-dimensional projector, so the encoder cost model,
batching law, burst-saturation result, and handoff bound carry over with no
caveat. On the language side a TP8 spot-check on the same shipped H100 data
shows the flagship decode-bound earlier: ITL p99 is already
\SI{87.7}{\milli\second} at the trace's own rate, so its knee sits at
$s{=}1$ (2.8\,req/s) against the 8B's $s{=}2$ --- decoder scale halves the
language pool's operating envelope while leaving the vision pool untouched,
widening the per-pool decoupling the thesis is about.

\Q{This is all vision encoders. Does anything generalize to audio, video, or any-to-any multimodal models?}
\A The mechanism is modality-agnostic; the constants are vision's. Nothing
in the argument uses that the encoder is a ViT --- it uses that a
disaggregated pre-stage has (a) a demand signal decoupled from request
rate, (b) heavy-tailed, long-range-dependent burstiness, and (c)
forecastability at provisioning horizons. Any modality whose work arrives
in bulk and is declared at admission --- audio segments, video frames, and
the frames class here is already video-like ($k^2$ grids up to $97^2$) ---
is a candidate for the same treatment, and any-to-any runtimes (vLLM-Omni)
multiply the pools and thus the signals to forecast, which is where
zero-shot forecasting earns its keep. I have measured one modality on one
trace; the method transfers, the $R^2$ and the knees do not.

\Q{Hardware moves fast --- faster cold starts, bigger models, new accelerators. Will this still matter in three years?}
\A The load-bearing parts are workload-structural, not hardware-specific,
and the trend runs \emph{toward} the thesis. The decoupling and the
long-range-dependent burstiness are properties of aggregated human demand,
continuous with thirty years of teletraffic measurement, not of any GPU.
Cold starts shrink but the encoder-pool lead is orchestration-bound (pod
scheduling, image pull), which weight-multicast advances like BlitzScale do
not touch. Bigger decoders only widen the gap: the 235B's language knee
sits an octave below the 8B's while the vision pool is unchanged. And the
calibration was built silicon-portable on purpose --- re-anchoring to a new
accelerator is one $f(N)$ benchmark, memoized in minutes, not a new thesis.

\Q{You make several ``to our knowledge, first'' claims. HeteroServe already costed a consumer-class encoder; Diao and Toner already benchmarked zero-shot forecasters on cloud demand. Are these firsts not narrow re-slicing?}
\A Each ``first'' is narrowed \emph{against exactly those works}, cited by
name, rather than around them. HeteroServe costed a consumer encoder for
cross-tier placement; my narrowed claim is first to integrate such a model
into a deployment \emph{planner} (aiconfigurator held no consumer-GPU
encoder data before this) and upstream it. Diao (Chronos on
Redshift-Serverless) and Toner are cited; the forecasting claim narrows to
the first \emph{provisioning-framed, rolling-origin} zero-shot benchmark on
the \emph{per-pool demand signals of a production multimodal trace},
consumed in closed loop. But I will not die defending adjective-laden
firsts: the thesis-level novelty is the \emph{inversion finding} --- the
least rate-visible signal is the most forecastable --- and I claim no single
ingredient as novel.
\Concede{The more qualifiers a ``first'' needs, the weaker it is. Retreat to
the finding, not the three narrow firsts; note the positioning table's
columns are the four ingredients the finding requires, valued at their
intersection.}

\Q{Decompose it: a one-trace characterization, a forecasting benchmark, an integration, a simulated loop. Which single piece is a novel research contribution rather than competent engineering? Is ``burstiest = most forecastable,'' one correlation on one week, big enough to carry a thesis?}
\A I concede no single component is a new \emph{method}, and I disclaim
ingredient novelty explicitly. The contribution is empirical: a falsifiable,
mechanism-grounded \emph{inversion} --- the signal most dangerous to reactive
scaling is the most forecastable --- that reorganizes how you provision,
demonstrated end to end (quantified, realized zero-shot, reproduced as
saturation behavior, closed in a loop, live-anchored on the control plane)
through a platform whose every seam is calibrated. If the bar is a novel
algorithm, I do not meet it and do not claim to; if it is a defended,
quantified, actionable finding that prior systems assumed the opposite of,
that is the thesis.
\Concede{The weight rests entirely on the finding's generality, untested
beyond one trace. Say that plainly rather than inflate any single
component.}

\Q{You did not build a system --- you glued AIC, Dynamo, and three off-the-shelf forecasters together. Where is the research?}
\A The research is not an artifact; it is a set of measured general claims
and the discipline that makes the gluing correct. The parts are
off-the-shelf on purpose --- forecasters consumed zero-shot, Dynamo's real
planner as the control path --- because the contribution is the
\emph{finding} (decoupled per-pool demand, the burstiest signal is the most
forecastable, mean-based conversion is the defect) and the \emph{composition}
that lets it be measured with per-pool attribution. And the places the
audit found errors were exactly the integration seams: measuring the wrong
execution engine, actuating on target rather than issue time, sizing a
denominator off an endpoint artifact. Getting those right ---
matched-engine, issue-time, ledgered --- is where the rigor lives, and the
by-products (the upstreamed dataset, the shipped-predictor comparison, the
actuation-semantics discipline) are independently usable. Gluing is easy;
gluing so the numbers survive ten adversarial rounds is the work.

\Q{Why is this doctoral-level work rather than excellent engineering?}
\A Engineering would be: pick the image metric, wire a forecaster into a
planner, show it works once. What makes this a thesis is that each link in
that chain is replaced by a \emph{measured general claim with its boundary
conditions}: a characterization result (the decoupling and its LRD structure,
continuous with the teletraffic literature), an impossibility-shaped result
(no averaging window rescues threshold-rule reactive control --- Fano
growth through 938 --- and the floorless matched-rating arm on the right
signal fails at every realizable lead; standing floors pass at higher
cost), a
feasibility result (the signal is forecastable where it must be, with
calibrated uncertainty), a mechanism result (the 1{,}616/1{,}717/736\,ms
triad separating anticipation from conversion --- the finding that mean-based
capacity conversion, not signal choice, is the defect), and a validated
platform others can re-anchor to new silicon with one benchmark. The
contribution survives its own ablations and ten adversarial audit rounds;
several individual pieces (the upstreamed dataset, the shipped-predictor
comparison, the actuation-semantics discipline) are independently usable.
Engineering delivers an artifact; this delivers the artifact plus the reasons
it must behave that way and the conditions under which it stops.

\Q{What would you do next, with another year?}
\A Three items, in priority order, each closing a boundary the thesis
explicitly draws. First, the one open validation item: a \emph{live encoder
worker} on the real request path --- every silicon anchor so far is a
measured-service replay (whole-tower forwards fed through the grid-faithful
FIFO), not a serving process; the manuscript names this as the remaining
Phase-E gap. Second, a \emph{second production trace}, the moment one with
per-request image counts exists: the single measurement that most needs
replication is the persistence result (Hurst ${\approx}0.93$, 69-minute
half-life), because the whole forecasting escape route rests on it. Third,
\emph{fine-tuned forecasters}: zero-shot was chosen as the conservative lower
bound, so trace-adapted models bound the headroom from above --- my
expectation, for what it is worth, is that they compress the p90 band and
move cost, not verdicts, because that is what the model-swap ablation already
showed. Further out: speculative encoder/prefill overlap at the seam (my
serialization is conservative by construction), and heterogeneous encoder
hardware, which the silicon-portable calibration was designed to make cheap
--- re-anchoring to a new accelerator is one $f(N)$ benchmark, not a new
thesis.

\Q{Is there anything you would now do differently from the start?}
\A Actuation semantics first. The most expensive lesson of the audit rounds
was that \emph{when} a forecast exists is as load-bearing as what it says ---
the target-time/issue-time distinction cost the headline a factor of two, and
I found it in review rather than in design. Done again, the issue-time
discipline (and both readiness models) would be in the harness from day one,
the way per-pool attribution was. Second, I would have measured the serving
engine's kernels \emph{before} building the HF-eager anchor, not after the
discrepancy forced it --- the matched-engine principle was learnable earlier.
Neither changes a conclusion; both would have saved months. And the ledger
habit --- every number bound to an artifact, a command, and the paragraph it
feeds --- I would keep unchanged; it is the reason ten adversarial rounds
tightened the thesis instead of dismantling it.

\Q{What did you learn that you did not expect?}
\A That the choice \emph{among the tested foundation models} barely matters.
I went in expecting model choice to be
a headline --- it is the axis with all the recent literature --- and within
the three tested TSFMs it turned
out to be second-order after capacity quantization: all three foundation
models pass the week at the identical \SI{651.1}{\milli\second} tail with only
cost shifting; swapping Dynamo's shipped predictors leaves verdicts unchanged;
integer-replica conversion quantizes away most differences \emph{between
them}. What
actually decides outcomes is the \emph{quantile} (p90 passes all 48 week-scale
cells on the encoder pool, p50 passes zero) and the \emph{conversion}
(tail-aware passes, mean-based fails --- with the same forecast). I will not
stretch that into ``information was never the scarce resource,'' because the
forward band's \emph{shape} is load-bearing: on burst3 the tested EWMA
alternatives split on outcome --- additive and width-scaled bands fail at p99
$1{,}285$ and $2{,}084$\,ms, and multiplicative EWMA passes only by carrying
substantially more standing work than the forecast arm. The correct statement
is narrower and still surprising: within the TSFM class the identity of the
model is second-order, while the state-adaptive upper edge, and what you do
with it, decide the verdict. I looked for a
forecasting result and found a control-theory result, and it is a better
finding than the one I wanted.

\Q{If the committee asks about philosophy-of-science framing (Lakatos, Popper).}
\A Only if raised --- I will not volunteer it. If it comes: the working
posture of the thesis is closest to Lakatos's methodology of research
programmes, honestly applied to my own work rather than to rivals. The hard
core is the per-pool provisioning claim; the cost model is protective belt;
and the test of integrity is what happens when the belt takes a hit. When the
$1.4$--$1.6\times$ silicon anomaly appeared, the repair was
\emph{content-increasing} in exactly Lakatos's sense: the execution-engine
hypothesis predicted something new and checkable --- that measuring the
serving engine's own kernels would close the gap --- and the dense 69-point
vLLM sweep then checked it, landing at 12.8\%. A progressive problemshift, not
an ad hoc rescue; the anomaly produced the thesis's strongest validation
artifact. Falsifiability I answer in Popper's currency because the committee
expects it (a trace with $R^2 \to 1$ kills the premise; a floorless
matched-rating reactive controller that holds the fast onsets at no more
than the forecast's cost kills the mechanism), while noting the honest
Lakatosian caveat that single results are rarely decisive --- which is why the
manuscript reports whole frontiers and both readiness models instead of point
verdicts. What I will \emph{not} do is use demarcation language against rival
approaches; that part of the apparatus did not survive (Laudan's 1983 verdict
on the demarcation problem is the standard citation), and a philosopher in the
room will know it.
\Trap{Keep it to ninety seconds. The framing is armor for the corrections
narrative --- ``the anomaly was repaired progressively, with a novel
prediction that was then verified'' --- not a lecture topic.}

% ==========================================================================
\section{Cram layer: rapid-fire numbers}
\label{sec:cram}

One-breath answers for facts I must produce instantly. If any of these takes
me more than five seconds in rehearsal, drill it again.

\begin{itemize}
\item \textbf{Trace:} 1M requests, exactly 7 days, Azure production
  multimodal cluster; 50\% text-only by construction; per-request image
  counts, no sizes.
\item \textbf{Frames class:} $k{>}10$ grids; 2.7\% of requests, 81.2\% of
  image volume; up to $97^2 = 9{,}409$ images in one request; 93--96\% of
  each burst's volume.
\item \textbf{Decoupling:} $R^2 = 0.123$ images-vs-requests ($0.654$ for
  decode offered-work); Jaccard 0.02 between top-percentile minutes;
  best-fit constant fraction 15.1 img/req.
\item \textbf{Burstiness:} image CV$^2$ 3.83 vs.\ request 0.32 at one
  minute; Fano 1.8 (1\,s) $\to$ 938 (1\,h); peak 370\,img/s ${\approx}16\times$
  weekly mean while requests stay under 4.5\,req/s.
\item \textbf{Persistence:} ACF half-life 69\,min (images) vs.\ 9\,min
  (requests); Hurst ${\approx}0.93$ vs.\ $0.83$ (spectral).
\item \textbf{Onsets:} 18 episodes; median ramp 6.5\,min, fastest 1\,min;
  median steepest single-minute jump 208\,img/s.
\item \textbf{Forecasting:} FMs win all 20 (signal, resolution) cells;
  images@1min 52.9\% WAPE (EWMA 54.2\%); images@1s 98.8\% (EWMA 146.1\%,
  no-skill 100\%); ${\approx}2$ WAPE
  points lost over 12 steps at 1-min bins; p90 coverage 0.88--0.90; p10--p90 band covers
  69\% vs.\ nominal 80\%.
\item \textbf{Cost model:} attention operator within ${\approx}3\%$ of
  physical; composed vs.\ vLLM tower within 12.8\% ($N{\le}64$, 69-point
  dense sweep); HF-eager gap 1.4--1.6$\times$ (engine, not model); unlocked
  clocks inflate ${\approx}17\%$; AIC SOL fallback ${\approx}5.9\times$
  optimistic; upstream PR \#1312.
\item \textbf{Knees:} order of magnitude apart on steady traffic (vision
  censored $\ge 32\times$, LLM first-over at $3\times$); 235B@TP8 knee one
  octave below 8B (ITL-bound, 87.7\,ms at trace rate).
\item \textbf{Closed loop:} 588\,ms zero-queueing floor on burst windows;
  calm day $\le$ 734\,ms; p90 passes 60/60 issue-time cells (both readiness
  models); p50 passes 23--24/60; the floorless matched-rating reactive arm
  loses burst3 at every lead (9.5--54\,s,
  never-scaled limit 54{,}170\,ms); reactive-with-headroom still fails it
  (1{,}973--11{,}568\,ms, despite up to $8\times$ tail cut); standing floors
  do pass but cost more (post-hoc common-warm $M{=}2$ $0.746\times$ static
  $=1.20\times$ forecast; shipped-mechanics $M{=}6$ $1.042\times$); triad on
  burst1@L60: 1{,}616 / 1{,}717 / 736\,ms.
\item \textbf{Economics:} $0.62\times$ = 2{,}410{,}755 / 3{,}869{,}278
  worker-seconds vs.\ static $C{=}8$ over 5.60 days; frontier $\rho10$--$20$
  $= 0.89$--$0.48\times$; strict re-warm 0.59--0.62$\times$; LLM week: fc-p90
  651.1\,ms at 518.6\,GPU-h vs.\ reactive fail 1{,}196.6\,ms at 281.8.
\item \textbf{Silicon closure:} winning schedule silicon-feasible, 588 $\to$
  784\,ms, ratio preserved; queue-replay 589 $\to$ 606\,ms at identical
  worker-seconds; static denominator 692\,ms, still feasible.
\item \textbf{Live anchor:} fc 695 vs.\ 651\,ms, reactive 3{,}531 vs.\
  3{,}495\,ms; TTFT within 1.1--6.8\%; replica-time within 8\%; one run per
  policy, mocked engines, fastest-onset window.
\item \textbf{Upstream status:} encoder dataset + system file are a
  \emph{draft} PR (ai-dynamo/aiconfigurator \#1312), \emph{not merged};
  data is commit-pinned (\texttt{1e8b9cac} locked-clock recollection,
  extensions through \texttt{b55dd0b7}) and snapshot-archived, so
  reproducibility does not depend on the merge.
\item \textbf{Cost units:} encoder in RTX-5090-class \emph{worker}-seconds
  (one worker $=$ one 5090-class GPU at TP1, 1:1), language in H100
  GPU-hours; never summed --- no single all-in system-cost number exists,
  only two per-class ratios ($0.62\times$ enc, $0.965\times$ lang).
\item \textbf{Batching law:} encoder cost depends only on total
  images/forward (no cross-image attention; verified to the microsecond
  across $(1,64)/(2,32)/(64,1)$); $f(1){=}8.16$, $f(64){=}316.4$\,ms;
  amortization saturates $N{\approx}6$--$8$; greedy FIFO, 64-image cap.
\item \textbf{Handoff:} $196{\times}4096$ bf16 ${\approx}1.6$\,MB/image
  (embedding-only), $\text{RTT}{+}\text{bytes}/\text{bw}$ on 25\,Gbps; burst
  p99 272\,ms (sets the 588 floor); DeepStack $4\times$
  ${\approx}6.4$\,MB/image needs ${\geq}100$\,Gbps; bandwidth-only, DMA not
  charged to the worker.
\item \textbf{SLA provenance:} TTFT/enc-delay p99 $<1$\,s, ITL p99
  $<100$\,ms $= 2\times$ Dynamo's 500/50\,ms window-mean defaults, applied
  as tails; joint ingress-to-first-token ${\approx}1.24$--$1.30$\,s at
  ${\approx}$p98 (encoder/prefill tails correlated $r{=}0.52$).
\item \textbf{Engine sweep} (vs.\ AIC): HF-SDPA 1.36--1.61$\times$ (rising
  with $N$); HF-FA2 1.25--1.47$\times$; \texttt{torch.compile}
  1.2--1.4$\times$; vLLM within 12.8\% for $N{\le}64$ ($+21.8\%$ only past
  the cap at $N{=}512$) --- vLLM being the engine AIC's kernels come from.
\item \textbf{235B flagship:} vision tower $+$ 4096 projector shared with 8B
  (encoder result transfers unchanged); TP8 language knee $s{=}1$
  (2.8\,req/s) vs 8B's $s{=}2$, ITL p99 87.7\,ms at trace rate --- decoder
  scale halves the language envelope.
\end{itemize}

\subsection*{Theory quick-facts card}
\addcontentsline{toc}{subsection}{Theory quick-facts card}

The same drill as above, for the definitional and theoretical facts. These
extend the list rather than repeat it: everything here is a number attached
to a \emph{term} the committee may probe, verified against the manuscript or
a tracked artifact.

\begin{itemize}
\item \textbf{SLA convention:} ``evaluation SLA'' $=$ the graded tail bounds
  (enc.\ delay p99, TTFT p99 $<\SI{1}{\second}$; ITL p99
  $<\SI{100}{\milli\second}$); ``SLO'' $=$ the general concept. Declared in
  the introduction, followed thesis-wide.
\item \textbf{ITL:} the per-gap equivalent of TPOT, the per-token decode gap
  TPOT summarizes; ITL $=$ TPOT in the live client.
\item \textbf{Percentile estimator:} order-statistic pick, no interpolation
  --- $\text{sorted}[\min(n{-}1,\lfloor np/100 \rfloor)]$, the converter's own
  rule, reused by the independent recompute. Every p99 is an observed delay,
  which is why plateau ties are exact (\SI{589.301}{\milli\second}).
\item \textbf{Headline in worker-hours:} 669.7 (forecast p90) against
  1{,}074.8 (deployable static $C{=}8$) $= 0.623\times$, over the 5.60-day,
  749{,}601-request span; \emph{no burst definition, SLA cut, or percentile
  selection enters either operand}.
\item \textbf{Threshold-integrated identity:}
  $\sum_i d_i = \int_0^\infty N(d_i>t)\,\mathrm{d}t$ --- total delay mass is
  the area under the full exceedance-count curve, so it integrates every
  threshold. Linear loss, declared population, post-hoc synthesis; it does
  not replace the every-window gate.
\item \textbf{Total delay mass (req$\cdot$s):} forecast \num{30277.5} vs
  reactive \num{34015.6} (11.0\% less) at \num{55.7} fewer worker-hours
  (7.7\%); static $C{=}8$ \num{25536.9}; warm floor $M{=}2$ \num{28343.3} at
  769.3 worker-hours. The floor interpolated to the forecast's delay burden
  costs 754.4 worker-hours, 12.6\% more.
\item \textbf{Excess mass} ($\sum(\text{delay}-\SI{1}{\second})^{+}$,
  req$\cdot$s): forecast \num{3688.0} vs reactive \num{6709.7} (45.0\% less);
  static $C{=}8$ \num{289.9}.
\item \textbf{Quantile ladder:} identical quantiles through p99
  (p50/p90/p95/p99, shared \SI{589.301}{\milli\second} plateau) while the CDFs
  already differ below it in violation counts; p99.5 752.8 vs 906.6; p99.99
  \num{6705.9} vs \SI{14724.0}{\milli\second}. The p99 tie is a statement
  about the plateau, not about the distributions.
\item \textbf{Cheapest passing reactive:} post-hoc common-warm $M{=}2$ costs
  $1.20\times$ the forecast at equal compliance on the symmetric warm-pool
  footing (equivalently the forecast uses 16.5\% less worker-time), and
  $1.67\times$ under the shipped strict re-warm mechanics ($M{=}6$). The
  $1.67\times$ is a sensitivity and \emph{must never be quoted alone}.
\item \textbf{Integer-millisecond crossings:} forecast's worst evaluation
  window \num{681.1}\,ms, so its every-window pass survives any bound down to
  \SI{682}{\milli\second} (warm-pool readiness); static $C{=}8$'s worst is
  \num{807.6}, so below \SI{808}{\milli\second} the static reference itself
  fails; the $M{=}6$ floor's is \num{906.6}, so below
  \SI{907}{\milli\second} the $1.042\times$ floor is no longer all-passing.
  Tightening the convention flips the comparators \emph{before} the forecast
  arm.
\item \textbf{Language week, below the contract statistic:} total TTFT mass
  \num{43598.8} vs reactive's \num{69671.3} request-seconds (37.4\% less);
  end-to-end mass 16.5\% less; p95 \num{178.9} vs \num{501.9}\,ms; ITL p99
  6.41 vs 21.28\,ms, both passing --- locating reactive's span-level failure
  in TTFT alone (651.1 vs \SI{1196.6}{\milli\second}).
\item \textbf{Issued lead arithmetic:} issue time $=$ target
  $-\,(h{-}1)\cdot\SI{15}{\second}$; $h{=}1$ carries \emph{zero} usable lead;
  an $H{=}12$ archive at \SI{15}{\second} bins supplies at most
  $(12{-}1)\times15 = \SI{165}{\second}$; declared grid
  $\{15,30,45,165\}$\,s. The old \SI{300}{\second} cell was target-indexed.
\item \textbf{Rolling origin:} $H{=}12$ bins, stride $H$ (non-overlapping,
  tiling); common span 167 origins at \SI{1}{\minute} up to \num{49727} at
  \SI{1}{\second}; index-thresholded, so 98.7\% of the week at
  \SI{1}{\second} but only the final \SI{33.4}{\hour} at \SI{1}{\minute}.
  Within-column comparisons only.
\item \textbf{WAPE:} $\sum_t|\hat y_t - y_t|/\sum_t y_t$; scale-free, defined
  on zero bins; no-skill floor 100\%. Pinball/CRPS: not computed.
\item \textbf{Paired-origin bootstrap:} 300 independently resampled
  time-ordered origins, not a day- or block-level population test. ARIMA
  66.8\% vs Chronos 66.0--66.1\% at \SI{15}{\second} is indistinguishable
  from zero; Toto+cov 60.9\% beats ARIMA with an interval excluding zero.
\item \textbf{Coverage, resolved:} p90 edge $0.876$--$0.903$ across all twelve
  horizon steps at \SI{15}{\second} bins and $0.857$--$0.911$ at
  \SI{1}{\minute}, no trend in $h$; $0.890$--$0.901$ in burst against
  $0.890$--$0.896$ outside; widest deviation over the twenty combinations
  $0.880$. Upper edge only.
\item \textbf{Persistence, all four signals:} ACF half-life images
  \SI{69}{\minute}, context \SI{15}{\minute}, requests \SI{9}{\minute},
  generated \SI{5}{\minute}; Hurst $0.93/0.86/0.83/0.91$. Half-life and
  Hurst orderings disagree, so no joint claim. Hurst is
  $H=(1{+}\beta)/2$ over 10\,min--12\,h; refit at 6\,h gives $0.90$, at 3\,h
  $0.82$.
\item \textbf{Overdispersion detail:} Fano $1.8$ (\SI{1}{\second}) $\to 938$
  (\SI{1}{\hour}); inter-arrival CV $1.55$; per-second counts overdispersed
  against Poisson$(1.65)$. Poisson would be Fano $=1$ at every scale.
\item \textbf{Rank coupling:} image count vs generation length, Spearman
  $-0.06$ at request level, no monotone trend in per-bucket median
  generation --- ``weakly rank-correlated,'' one of three tested
  relationships.
\item \textbf{The 29\% mechanism:} backlog above 1{,}000 queued images for
  805\,s across the 8.3-hour window, peaking at 4{,}701, worst second landing
  2{,}726 images; encoder-delay p99 \SI{5.0}{\second} ($5\times$ the bound);
  1-minute busy fraction reaches $0.90$ with 0.6\% of minutes above an $0.8$
  target, while the minute-mean reads healthy. Deterministic FIFO replay,
  \emph{not} an analytic queueing model.
\item \textbf{Knee detail:} vision uncensored through $s{=}32$
  ($\geq 256$\,img/s); language $s{=}2$ calibrated (5.5\,req/s) or $s{=}3$
  polynomial (8.3\,req/s), both ITL-bound with $1.8$--$2.1\times$ TTFT slack
  remaining; at the calibrated language knee the encoder idles at 4.5\%
  utilization; on the burst window neither pool has a knee (language
  \SI{1868}{\milli\second}, vision \SI{5013}{\milli\second} at $s{=}1$).
\item \textbf{ViT arithmetic:} patch 16, $448^2 \to 28\times28 = 784$
  patches; $2\times2$ PatchMerger $\to 196$ post-merge tokens/image;
  $196\times4096$ bf16 ${\approx}\SI{1.6}{\mega\byte}$ shipped per image. The
  trace's \texttt{ContextTokens} instead carry ${\approx}66$ tokens/image,
  roughly $3\times$ below the Qwen3-VL-consistent convention.
\item \textbf{Tower:} 27 blocks, hidden 1152, 16 heads $\times$ 72,
  GELU MLP 4304, PatchMerger into 4096, three DeepStack mergers; shared
  across the 8B/32B/235B family.
\item \textbf{Parallelism layout:} TP2$\times$2 $=$ two tensor-parallel-2
  replica groups (four H100s, iso-cost with TP1$\times$4). TP2$\times$2
  passes all five burst windows at 763--\SI{846}{\milli\second}, a
  $1.55$--$1.74\times$ tail cut over TP1$\times$4, which fails every one; at
  TP1 no replica count is feasible (p99 pinned near $1{,}285$\,ms even at
  $N{=}16$).
\item \textbf{Cold-start anatomy:} encoder model-side minimum
  \SIrange{3}{5}{\second}, above ${\approx}\SI{5}{\second}$
  orchestration-bound (pod scheduling, image pull, node acquisition);
  language lead is the weight-loading regime, \SIrange{60}{300}{\second}.
  Encoder comparator sweep $L\in\{15,30,45,300\}$\,s; forecast arm
  $\{15,30,45,165\}$\,s.
\item \textbf{Billing instants:} language GPU-hours bill from instance
  \emph{order} (so no cold-start surcharge may be added to them), the encoder
  pool from \emph{readiness} with cold start reported separately as a warm-up
  column. A second reason the classes are never summed.
\item \textbf{Masking, measured:} the reactive arm's aggregate p99 clears the
  bound while it fails the fastest onset, whose ${\approx}3{,}100$ requests
  are under half a percent of the span. Hence aggregate \emph{and}
  per-window, jointly; the per-window grid doubles as the masking check.
\item \textbf{Finer-grain audit:} across the seven daily folds the pure
  forecast fails two --- F6 at \SI{1264}{\milli\second} (a 7-minute onset,
  1-minute peak, \SI{263}{img/s} needing ${\approx}18$ workers at $\rho{=}15$
  where the band places 16), recovered post hoc by the hybrid combiner at
  \SI{679}{\milli\second} and $0.76\times$ static aggregate ($0.79\times$ on
  F6); and F2 at \SI{1319}{\milli\second}, which is an
  \emph{archive-coverage} failure, not an onset miss.
\item \textbf{Readiness definitions:} warm-pool $=$ a scaled-down replica
  returns without re-paying cold start; strict re-warm $=$ it pays the lead
  again. Shipped mechanics are asymmetric (forecast/oracle warm, reactive
  strict); both directions repaired and measured; the forecast's strict
  variant retains matured orders as reusable for a \SI{165}{\second} hold, so
  that endpoint is a sensitivity, not a fully equalized model.
\end{itemize}
% ==========================================================================
\section{Defense mechanics and delivery}

A different kind of preparation: not what to say, but how to say it in the
room. Rehearse these like the answers above.

\Q{Draw your system on the board.}
\A Five boxes, two arrows, in this order. (1)~The trace --- 1M requests,
7 days, each row carrying an image count and token counts. (2)~The per-pool
provisioner above it --- image rate $\to$ vision workers by \emph{forecast};
token rates $\to$ language replicas by tail-aware \emph{conversion}. (3)~The
vision pool --- $C(t)$ RTX~5090 encoder workers, a FIFO costed by AIC's
$f(N)$, SLA encoder-delay p99 $<\SI{1}{\second}$. (4)~The seam --- each
request arrives at the language pool its encoder delay later. (5)~The
language pool --- $N\times$TP H100 workers, \texttt{dynamo.replay} with a
KV router, SLA TTFT p99 $<\SI{1}{\second}$, ITL p99 $<\SI{100}{\milli\second}$.
Draw the dashed demand arrow from the trace up to the provisioner last ---
that is the whole thesis: the vision arrow does not track the request arrow.
\Point{This is the architecture figure. If I can draw it from memory in
thirty seconds, the framing is set before the questions start.}

\Q{The ninety-second opening, before any question.}
\A ``Multimodal serving runs the vision encoder in its own pool, scaled
separately from the language model. The question is what signal sizes that
pool. On the first public production multimodal trace, the image demand that
drives the encoder is almost decoupled from the request rate every
autoscaler watches --- $R^2{=}0.12$ --- an order of magnitude burstier, and
its peaks do not even coincide. So the encoder needs its own signal. Watching it with a floorless,
matched-rating trailing controller is not enough: bursts peak in minutes
against cold starts of minutes, and that arm loses the fastest onset at
every issued lead (standing reactive floors do pass, at higher cost). The escape is that the same signal is the most \emph{persistent} of
the four, so a pretrained forecaster reads it zero-shot --- and provisioning
on that forecast holds the tail SLA on every declared evaluation window,
including the fast onset the matched-rating reactive controller loses, at
0.62 times the static fleet, and more cheaply than every tested
evaluation-passing floor in the reported sweep.'' Then stop.
\Trap{Do not front-load the $0.62\times$. Open on the decoupling; the
economics are the payoff, not the hook.}

\Q{A committee member says: ``I don't follow your setup.''}
\A Drop the jargon and use the analogy. ``A multimodal request is part
pictures, part text. We process the pictures on one set of machines and the
text on another, and we choose how many of each. Everyone sizes both by
counting \emph{requests}. But one request can carry one picture or nine
thousand, so counting requests tells you almost nothing about the picture
machines --- you have to count \emph{pictures}, and the picture count is
spiky in a way request count is not. My thesis: count pictures, and
\emph{predict} the count, because by the time you see the spike it is too
late to boot machines.'' Then offer to go as deep as they want.

\Q{If time is short, which single result do I show?}
\A The burst3 reactive-versus-forecast split: same signal, same lead, same
platform --- the floorless matched-rating reactive arm fails the fast onset
(p99 up to \SI{54}{\second}), the
p90 forecast holds the \SI{588}{\milli\second} floor. It carries the whole
argument in one comparison. If they want the \emph{why}, the
1{,}616/1{,}717/\SI{736}{\milli\second} triad; if they want the premise, the
$R^2{=}0.123$ decoupling. Everything else is support.

\Q{A question I genuinely cannot answer.}
\A The protocol, decided in advance: say ``I don't know'' or ``I did not
test that'' in the first sentence, then locate the boundary --- which side
of measured/modeled it falls on, and why I stopped where I did --- and offer
the experiment that would answer it. I do not bluff a number; the one thing
that cannot be recovered in the room is credibility. If it is a real gap,
conceding it fast is cheaper than defending it badly.

\Q{An examiner challenges a specific number and I cannot reconstruct it live.}
\A Point to the ledger, without defensiveness. Every number in the document
is generated by a versioned script over archived inputs, and each has a
results-ledger entry naming the artifact, the exact command, the date, and
the paragraph it feeds. ``That figure comes from \texttt{[artifact]}; I can
show the command and its provenance --- the ledger and the reproduction
manifest exist so no number is unaccountable.'' The discipline is the
answer.

\Q{The hostile examiner.}
\A Tone is the whole game, and it is mine to set, not theirs. Slow down,
agree with the part that is right, concede the real limitation in one
sentence, and separate it from what survives --- without heat. The audit
history is armor: every hard number was already moved against my own
interest in review, so there is nothing to be defensive about. If a question
is an objection wearing a question's clothes, name the objection plainly and
answer \emph{that}. Confidence reads as having-audited-it; defensiveness
reads as fearing-they-will.

\end{document}
